\documentclass[../main.tex]{subfiles}
\begin{document}

\chapter{スレッド性能上の考慮事項}
\lessoninfo{
  video={P2L5},
  textbook={OSTEP \cite{ostep} ch.~26，33；Silberschatz ら \cite{silberschatz2018} ch.~4--5；Kerrisk \cite{kerrisk2010} ch.~49--50，63；Pai ら \cite{pai1999}},
  keywords={性能指標，実行時間，スループット，応答時間，マルチプロセスモデル，マルチスレッドモデル，イベント駆動モデル，イベントディスパッチャ，非同期I/O，ヘルパ，AMPED，Flash，ワークロード，ベースライン},
}

第3章から第5章ではスレッド，そのプログラミングインタフェース，およびその
実装を導入し，マルチスレッド化がアプリケーションをより速く，より効率的に
すると論じた．この章では，そのような主張をどのように評価するのかを問う．
どの設計が優れているかは性能指標とワークロードに依存することを示し，ウェブ
サーバのマルチプロセス，マルチスレッド，イベント駆動の各アーキテクチャを
比較し，実験の設計の事例研究として Pai ら \cite{pai1999} の Flash ウェブ
サーバの評価をたどる．

\section{スレッドモデルの比較}
\label{sec:l06-compare}

第3章では，ボス・ワーカーパターンとパイプラインパターンを，一群の要求を
処理し終えるのに要する総時間で比較した．これは妥当な基準のうちの1つに
すぎず，別の基準を取れば判定は逆転しうる．\source{P2L5，どのスレッド
モデルが優れているか．}5つのスレッドを持つサーバが $n$ 個の要求の一群を
受け取り，各要求の処理に \SI{125}{\milli\second} を要するとする．
\begin{itemize}
  \item ボス・ワーカー構成では，1つのスレッドがボスであり，その要求当たりの
    時間は無視でき，$w = 4$ 個のスレッドがワーカーで，そのそれぞれが
    $t = \SI{125}{\milli\second}$ で1つの要求全体を処理する．
  \item パイプライン構成では，$k = 5$ 個のスレッドのそれぞれがステージで
    あり，どのステージも $s = \SI{25}{\milli\second}$ を要する．
\end{itemize}
ボス・ワーカーでは要求は $w$ 個ずつ完了するので，要求 $i$ は時刻
$\lceil i/w \rceil\, t$ に完了する．パイプラインでは要求 $i$ は
$(k + i - 1)\, s$ に完了する．これらの完了時刻から2つの指標が導かれる．%
\caution{$n$ 個の要求はすべて時刻 0 に存在しており，完了はそこから数えて
いる．到着が時間的に散らばる場合に意味を持つ指標は応答時間
（\cref{sec:l06-metrics}）である．}%
一群全体を処理し終える時間は，最後の要求の完了時刻である．1つの要求の
\emph{平均完了時間}は
\begin{equation}
  \overline{T}_{\mathrm{bw}}(n) = \frac{t}{n} \sum_{i=1}^{n}
    \Bigl\lceil \frac{i}{w} \Bigr\rceil ,
  \qquad
  \overline{T}_{\mathrm{pl}}(n) = k\, s + \frac{(n-1)\, s}{2}
  \label{eq:l06-mean-completion}
\end{equation}
である．

\begin{figure}[htbp]
  \centering
  % Completion times of nine requests: boss-workers with four workers
% (100 ms per request) and a five-stage pipeline (20 ms per stage).
% Usage: % Completion times of nine requests: boss-workers with four workers
% (100 ms per request) and a five-stage pipeline (20 ms per stage).
% Usage: % Completion times of nine requests: boss-workers with four workers
% (100 ms per request) and a five-stage pipeline (20 ms per stage).
% Usage: \input{figures/l06-completion}   (width ~116 mm)
\begin{tikzpicture}[x=1mm, y=1mm, font=\footnotesize\sffamily,
  >={Stealth[length=1.6mm]},
  lab/.style={font=\scriptsize\sffamily},
  row/.style={font=\scriptsize\sffamily, anchor=east},
  done/.style={draw, circle, fill=ln-accent!20, minimum size=4.2mm,
    inner sep=0pt, font=\scriptsize\sffamily},
  mean/.style={ln-caution, densely dashed, thick}]
  % x = 24 + 0.3 t, with t in ms
  % boss-workers
  \node[row] at (22,18) {\iflangja{ボス・ワーカー}{boss-workers}};
  \draw[ln-rule] (24,18) -- (114,18);
  \draw[mean] (74,15) -- (74,23);
  \node[lab, anchor=south, text=ln-caution] at (74,23.1)
    {\iflangja{平均 167\,ms}{mean 167\,ms}};
  \node[done] at (54,18) {4};
  \node[done] at (84,18) {4};
  \node[done] at (114,18) {1};
  % pipeline
  \node[row] at (22,4) {\iflangja{パイプライン}{pipeline}};
  \draw[ln-rule] (24,4) -- (114,4);
  \draw[mean] (78,-0.5) -- (78,8.5);
  \node[lab, anchor=south, text=ln-caution] at (78,8.6)
    {\iflangja{平均 180\,ms}{mean 180\,ms}};
  \foreach \t in {100,120,...,260} { \node[done] at (24+0.3*\t,4) {1}; }
  % time axis
  \draw[->] (24,-4) -- (116,-4);
  \foreach \t in {0,50,...,300} {
    \draw (24+0.3*\t,-4) -- (24+0.3*\t,-5);
    \node[lab, anchor=north] at (24+0.3*\t,-5.2) {\t};
  }
  \node[row] at (22,-4) {\iflangja{時間 (ms)}{time (ms)}};
\end{tikzpicture}
   (width ~116 mm)
\begin{tikzpicture}[x=1mm, y=1mm, font=\footnotesize\sffamily,
  >={Stealth[length=1.6mm]},
  lab/.style={font=\scriptsize\sffamily},
  row/.style={font=\scriptsize\sffamily, anchor=east},
  done/.style={draw, circle, fill=ln-accent!20, minimum size=4.2mm,
    inner sep=0pt, font=\scriptsize\sffamily},
  mean/.style={ln-caution, densely dashed, thick}]
  % x = 24 + 0.3 t, with t in ms
  % boss-workers
  \node[row] at (22,18) {\iflangja{ボス・ワーカー}{boss-workers}};
  \draw[ln-rule] (24,18) -- (114,18);
  \draw[mean] (74,15) -- (74,23);
  \node[lab, anchor=south, text=ln-caution] at (74,23.1)
    {\iflangja{平均 167\,ms}{mean 167\,ms}};
  \node[done] at (54,18) {4};
  \node[done] at (84,18) {4};
  \node[done] at (114,18) {1};
  % pipeline
  \node[row] at (22,4) {\iflangja{パイプライン}{pipeline}};
  \draw[ln-rule] (24,4) -- (114,4);
  \draw[mean] (78,-0.5) -- (78,8.5);
  \node[lab, anchor=south, text=ln-caution] at (78,8.6)
    {\iflangja{平均 180\,ms}{mean 180\,ms}};
  \foreach \t in {100,120,...,260} { \node[done] at (24+0.3*\t,4) {1}; }
  % time axis
  \draw[->] (24,-4) -- (116,-4);
  \foreach \t in {0,50,...,300} {
    \draw (24+0.3*\t,-4) -- (24+0.3*\t,-5);
    \node[lab, anchor=north] at (24+0.3*\t,-5.2) {\t};
  }
  \node[row] at (22,-4) {\iflangja{時間 (ms)}{time (ms)}};
\end{tikzpicture}
   (width ~116 mm)
\begin{tikzpicture}[x=1mm, y=1mm, font=\footnotesize\sffamily,
  >={Stealth[length=1.6mm]},
  lab/.style={font=\scriptsize\sffamily},
  row/.style={font=\scriptsize\sffamily, anchor=east},
  done/.style={draw, circle, fill=ln-accent!20, minimum size=4.2mm,
    inner sep=0pt, font=\scriptsize\sffamily},
  mean/.style={ln-caution, densely dashed, thick}]
  % x = 24 + 0.3 t, with t in ms
  % boss-workers
  \node[row] at (22,18) {\iflangja{ボス・ワーカー}{boss-workers}};
  \draw[ln-rule] (24,18) -- (114,18);
  \draw[mean] (74,15) -- (74,23);
  \node[lab, anchor=south, text=ln-caution] at (74,23.1)
    {\iflangja{平均 167\,ms}{mean 167\,ms}};
  \node[done] at (54,18) {4};
  \node[done] at (84,18) {4};
  \node[done] at (114,18) {1};
  % pipeline
  \node[row] at (22,4) {\iflangja{パイプライン}{pipeline}};
  \draw[ln-rule] (24,4) -- (114,4);
  \draw[mean] (78,-0.5) -- (78,8.5);
  \node[lab, anchor=south, text=ln-caution] at (78,8.6)
    {\iflangja{平均 180\,ms}{mean 180\,ms}};
  \foreach \t in {100,120,...,260} { \node[done] at (24+0.3*\t,4) {1}; }
  % time axis
  \draw[->] (24,-4) -- (116,-4);
  \foreach \t in {0,50,...,300} {
    \draw (24+0.3*\t,-4) -- (24+0.3*\t,-5);
    \node[lab, anchor=north] at (24+0.3*\t,-5.2) {\t};
  }
  \node[row] at (22,-4) {\iflangja{時間 (ms)}{time (ms)}};
\end{tikzpicture}

  \caption{9個の要求の完了時刻．各円はその時刻に完了する要求の数を示す．
    4つのワーカーを持つボス・ワーカーは \SI{125}{\milli\second} ごとに
    4個の要求を完了し，5段のパイプラインは最初の \SI{125}{\milli\second}
    の後は \SI{25}{\milli\second} ごとに1個の要求を完了する．破線は平均
    完了時間を示す．}
  \label{fig:l06-completion}
\end{figure}

\begin{example}
  \label{ex:l06-completion}
  $n = 9$ 個の要求の一群（\cref{fig:l06-completion}）では，ボス・ワーカーは
  \SI{125}{\milli\second} に4個，\SI{250}{\milli\second} に4個，最後の1個を
  \SI{375}{\milli\second} に完了する．パイプラインは最初の要求を
  \SI{125}{\milli\second} に完了し，その後は \SI{25}{\milli\second} ごとに
  1個ずつ，最後の1個を \SI{325}{\milli\second} に完了する．一群を処理し
  終えるのはパイプラインの方が早い．しかし \cref{eq:l06-mean-completion}
  によれば，平均完了時間はボス・ワーカーが
  $(4 \cdot 125 + 4 \cdot 250 + 375)/9 \approx \SI{208.3}{\milli\second}$，
  パイプラインが $125 + 8 \cdot 25/2 = \SI{225}{\milli\second}$ であり，
  平均的な要求に応じるのはボス・ワーカーの方が早い．$n = 40$ では両方の
  指標でパイプラインが勝つ．一群を処理し終えるのが
  \SI{1250}{\milli\second} ではなく \SI{1100}{\milli\second}，平均が
  \SI{687.5}{\milli\second} ではなく \SI{612.5}{\milli\second} である．
\end{example}

ボス・ワーカーは $t$ の後に $w$ 個の要求をまとめて完了するので，多くの要求
が早く終わる．一方パイプラインは1ステージ時間当たり1個の要求しか送り出さ
ない．しかし定常状態では，パイプラインが毎秒 $1/s = 40$ 個の要求を完了
するのに対しボス・ワーカーは $w/t = 32$ 個であり，大きな一群ではこの高い
完了率が支配的になる．したがってどちらのモデルが優れているかは，指標にも
要求の数にも依存する．

\needspace{6\baselineskip}
\section{スレッドの有用性}
\label{sec:l06-useful}

これまでの章では，スレッドを使う理由がいくつか挙げられた．
\source{P2L5，スレッドは有用か，視覚的な比喩；OSTEP ch.~26．}
\begin{itemize}
  \item \textbf{並列化}：複数の CPU 上では，同じタスクの一部分ずつを担当
    するスレッドがそのタスクを高速化する．
  \item \textbf{特化}：特定のタスクに専念するスレッドはホットキャッシュで
    動作し，異なる優先度を与えることもできる．
  \item \textbf{効率}：スレッドはアドレス空間を共有するので，プロセスより
    必要なメモリが少なく，通信と同期も安価である．
  \item \textbf{アイドル時間の隠蔽}：単一の CPU 上でも，あるスレッドが
    I/O を待つ間に別のスレッドが実行できる．
\end{itemize}
これらの利点がスレッドに基づく設計を有用にするかどうかは，何を測るかに
依存し，立場が異なれば測るものも異なる．行列を乗算するアプリケーションに
とって意味のある尺度は，その実行時間である．ウェブサービスにとっては，
単位時間当たりに処理するクライアントの要求数と，クライアントが応答を待つ
時間であり，後者は平均，最大，あるいは 95 パーセンタイルのような
パーセンタイルとして報告されうる．ハードウェアの所有者にとっては，CPU など
の資源の使用率である．したがって，要求に応じるあらゆるシステムについて，
3つの尺度が繰り返し現れる．単位時間当たりに完了する要求の数，1つの要求に
応答するまでの時間，そしてシステムの能力のうちどれだけが使われているか，
である．

\begin{example}
  あるサーバが 200 個の要求を処理し，そのうち 180 個は
  \SI{20}{\milli\second} 後に，20 個は \SI{400}{\milli\second} 後に応答を
  受け取るとする．平均応答時間は
  $(180 \cdot 20 + 20 \cdot 400)/200 = \SI{58}{\milli\second}$ である．
  95 パーセンタイルとは，少なくとも 95\,\% の要求がそれを超えない最小の
  値であり，昇順に並べたときの 190 番目の応答時間は
  \SI{400}{\milli\second} なので，95 パーセンタイルは
  \SI{400}{\milli\second} である．平均は，10 個に1個の要求が大多数の 20 倍
  待つという事実を覆い隠す．
\end{example}

\needspace{6\baselineskip}
\section{性能指標}
\label{sec:l06-metrics}

ある設計が別の設計より良い性能を示すという主張は，指標を定めて初めて意味を
持つ．\source{P2L5，性能指標，スレッドは本当に有用か；Silberschatz ら
ch.~5．}

\begin{definition}[性能指標]
  システムの測定基準，すなわち対象とするシステムの，測定可能で定量化できる
  性質であって，その振る舞いを評価する---例えば別の実装と比較する---ために
  使えるものを\term[せいのうしひょう]{性能指標}[performance metric]という．
\end{definition}

指標は理想的には，実際のマシンに配備された実際のソフトウェアを実際の
ワークロードの下で実験して得られる．これはしばしば実行不可能である．マシン
が利用できないかもしれず，ソフトウェアが完成していないかもしれず，実際の
利用者に実験のために待つよう求めることもできない．そこで測定は代わりに，
現実的な状況を代表するように設計された小規模な実験や，シミュレーションから
得られる．このように，実験を行って指標を得るために用意された環境を
\term[てすとべっど]{テストベッド}[testbed]という．

最もよく使われる指標は次のものである．
\begin{description}
  \item[実行時間] タスクの\term[じっこうじかん]{実行時間}[execution time]
    とは，その開始から完了までの時間である．行列乗算のようなバッチ計算に
    とって自然な指標である．
  \item[スループット] システムの\term[するーぷっと]{スループット}%
    [throughput]とは，単位時間当たりに完了するタスク，例えば要求の数で
    ある．これは\emph{要求レート}，すなわち単位時間当たりに到着する要求の
    数とは異なる．要求レートに追いつけないシステムでも，要求を完了する
    速さはそのスループットにとどまり，未処理の要求は溜まっていく．
  \item[応答時間] 要求の\term[おうとうじかん]{応答時間}[response time]と
    は，その発行からクライアントが応答を受け取るまでの時間であり，処理
    される前に要求が待つ時間も含む．多数の要求については，平均，最大，
    あるいはパーセンタイルで要約される．
  \item[待ち時間] タスクの\term[まちじかん]{待ち時間}[wait time]とは，その
    投入から実行が始まるまでの時間である．利用者が，ジョブがいつ終わるか
    よりいつ始まるかを気にする場合に重要になる．
  \item[CPU 使用率] \term[CPUしようりつ]{CPU使用率}[CPU utilization]とは，
    CPU が有用な仕事をしている時間の割合である．他の資源の使用率も同じ
    ように定義される．
\end{description}

さらに多くの指標が，性能をその費用や特定の立場の関心に結び付ける
（\cref{tab:l06-metrics}）．サービスの提供者はしばしば，
\term{SLA}[service level agreement]%
\index{service level agreement|see{SLA}}（サービスレベル合意）で性能目標を
約束する．例えば，要求の 99\,\% が \SI{200}{\milli\second} 以内に応答を
受け取る，といったものである．そのとき，合意に違反した要求の割合がそれ自体
1つの指標となる．

\begin{table}[htbp]
  \centering
  \caption{性能指標と，典型的にそれに関心を持つ立場．}
  \label{tab:l06-metrics}
  \small
  \begin{tabular}{@{}>{\raggedright\arraybackslash}p{0.25\linewidth}>{\raggedright\arraybackslash}p{0.45\linewidth}>{\raggedright\arraybackslash}p{0.22\linewidth}@{}}
    \toprule
    指標 & 測るもの & 関心を持つ立場 \\
    \midrule
    実行時間 & タスクを完了するまでの時間 & バッチジョブの利用者 \\
    スループット，要求レート & 単位時間当たりに完了するタスク，到着する
      要求 & 運用者 \\
    応答時間 & 要求に応答するまでの時間 & クライアント \\
    待ち時間 & タスクが開始するまでの時間 & 利用者 \\
    CPU 使用率 & CPU が動いている時間の割合 & 運用者 \\
    プラットフォーム効率 & 使用した資源に対するスループット &
      データセンタの運用者 \\
    ドル当たり・ワット当たり性能 & 費用や電力に対する性能 & 運用者，購入者 \\
    SLA 違反 & 合意した目標を外れた要求の割合 & 提供者，クライアント \\
    クライアント体感性能 & クライアントから見た品質．例えば動画のフレーム
      レート & クライアント \\
    総合性能 & 多数のタスクやクライアントにわたって平均した指標 & 運用者 \\
    平均資源使用量 & 例えば時間にわたって使用されたメモリ & 運用者 \\
    \bottomrule
  \end{tabular}
\end{table}

\subsection{場合による}

したがって，スレッドが有用かという問いに単一の答えはない．答えは指標に
依存する．\cref{sec:l06-compare}が平均完了時間と総完了時間について示した
とおりである．そして答えはワークロードにも依存する．2つのスレッドモデルの
どちらが優れているかは，要求の数とともに変わった．同じことが，オペレー
ティングシステムのほとんどの設計上の判断について成り立つ．グラフの種類が
違えば有利な最短経路アルゴリズムも違い，ファイルアクセスの型が違えば有利な
ファイルシステムの設計も違い，要求の負荷が違えば有利なサーバの
アーキテクチャも違う．

\begin{keypoint}
  一般に最良の設計というものはない．比較が意味を持つのは指標とワークロード
  を明示した場合に限られ，「どちらが優れているか」という問いへの答えは，
  たいてい「場合による」である．
\end{keypoint}

\needspace{6\baselineskip}
\section{マルチプロセスサーバとマルチスレッドサーバ}
\label{sec:l06-mp-mt}

この章の残りでは，並行サーバを構成する方法を，静的なファイルを提供する
ウェブサーバを例に比較する．\source{P2L5，マルチプロセス対マルチスレッド；
Silberschatz ら ch.~4．}クライアントが要求を送り，サーバは
\cref{fig:l06-web-steps}の各ステップでそれを処理し，ファイルを送って応答
する．要求の解析や応答ヘッダの計算のように，CPU 時間だけを必要とする
ステップもある．他のステップはブロックしうる．接続の受け付けはクライアント
を待ち，要求の読み込みとデータの送信はネットワークを待ち，ファイルの検索と
読み込みはディスクを待つ場合がある．

\begin{figure}[htbp]
  \centering
  % Processing steps of a simple web server for static content; the bars
% mark the steps that may wait for the network or for the disk.
% Usage: % Processing steps of a simple web server for static content; the bars
% mark the steps that may wait for the network or for the disk.
% Usage: % Processing steps of a simple web server for static content; the bars
% mark the steps that may wait for the network or for the disk.
% Usage: \input{figures/l06-web-steps}   (width ~116 mm)
\begin{tikzpicture}[x=1mm, y=1mm, font=\scriptsize\sffamily,
  st/.style={draw, rounded corners=1.5pt, fill=ln-box, align=center,
    minimum width=14.6mm, minimum height=9mm, inner sep=0.5pt},
  num/.style={font=\tiny\sffamily, text=ln-muted, anchor=south},
  net/.style={draw=none, fill=ln-accent!60},
  dsk/.style={draw=none, fill=ln-caution!60},
  >={Stealth[length=1.2mm]}]
  \node[st] (s1) at (7.3,0)   {\iflangja{接続の\\受け付け}{accept\\connection}};
  \node[st] (s2) at (24.2,0)  {\iflangja{要求の\\読み込み}{read\\request}};
  \node[st] (s3) at (41.1,0)  {\iflangja{要求の\\解析}{parse\\request}};
  \node[st] (s4) at (58.0,0)  {\iflangja{ファイル\\の検索}{find\\file}};
  \node[st] (s5) at (74.9,0)  {\iflangja{ヘッダ\\の計算}{compute\\header}};
  \node[st] (s6) at (91.8,0)  {\iflangja{ヘッダ\\の送信}{send\\header}};
  \node[st] (s7) at (108.7,0) {\iflangja{ファイル\\読込・送信}{read file,\\send data}};
  \foreach \i [evaluate=\i as \j using int(\i+1)] in {1,...,6} {
    \draw[->] (s\i.east) -- (s\j.west);
  }
  \foreach \i in {1,...,7} { \node[num] at (s\i.north) {\i}; }
  % may wait for the network
  \foreach \x in {7.3,24.2,91.8,108.7} {
    \fill[net] (\x-7.3,-6.2) rectangle (\x+7.3,-5.0);
  }
  % may wait for the disk
  \foreach \x in {58.0,108.7} {
    \fill[dsk] (\x-7.3,-8.2) rectangle (\x+7.3,-7.0);
  }
  % legend
  \fill[net] (22,-12.6) rectangle (27,-11.4);
  \node[anchor=west] at (28,-12) {\iflangja{ネットワークを待つ可能性}{may wait for the network}};
  \fill[dsk] (68,-12.6) rectangle (73,-11.4);
  \node[anchor=west] at (74,-12) {\iflangja{ディスクを待つ可能性}{may wait for the disk}};
\end{tikzpicture}
   (width ~116 mm)
\begin{tikzpicture}[x=1mm, y=1mm, font=\scriptsize\sffamily,
  st/.style={draw, rounded corners=1.5pt, fill=ln-box, align=center,
    minimum width=14.6mm, minimum height=9mm, inner sep=0.5pt},
  num/.style={font=\tiny\sffamily, text=ln-muted, anchor=south},
  net/.style={draw=none, fill=ln-accent!60},
  dsk/.style={draw=none, fill=ln-caution!60},
  >={Stealth[length=1.2mm]}]
  \node[st] (s1) at (7.3,0)   {\iflangja{接続の\\受け付け}{accept\\connection}};
  \node[st] (s2) at (24.2,0)  {\iflangja{要求の\\読み込み}{read\\request}};
  \node[st] (s3) at (41.1,0)  {\iflangja{要求の\\解析}{parse\\request}};
  \node[st] (s4) at (58.0,0)  {\iflangja{ファイル\\の検索}{find\\file}};
  \node[st] (s5) at (74.9,0)  {\iflangja{ヘッダ\\の計算}{compute\\header}};
  \node[st] (s6) at (91.8,0)  {\iflangja{ヘッダ\\の送信}{send\\header}};
  \node[st] (s7) at (108.7,0) {\iflangja{ファイル\\読込・送信}{read file,\\send data}};
  \foreach \i [evaluate=\i as \j using int(\i+1)] in {1,...,6} {
    \draw[->] (s\i.east) -- (s\j.west);
  }
  \foreach \i in {1,...,7} { \node[num] at (s\i.north) {\i}; }
  % may wait for the network
  \foreach \x in {7.3,24.2,91.8,108.7} {
    \fill[net] (\x-7.3,-6.2) rectangle (\x+7.3,-5.0);
  }
  % may wait for the disk
  \foreach \x in {58.0,108.7} {
    \fill[dsk] (\x-7.3,-8.2) rectangle (\x+7.3,-7.0);
  }
  % legend
  \fill[net] (22,-12.6) rectangle (27,-11.4);
  \node[anchor=west] at (28,-12) {\iflangja{ネットワークを待つ可能性}{may wait for the network}};
  \fill[dsk] (68,-12.6) rectangle (73,-11.4);
  \node[anchor=west] at (74,-12) {\iflangja{ディスクを待つ可能性}{may wait for the disk}};
\end{tikzpicture}
   (width ~116 mm)
\begin{tikzpicture}[x=1mm, y=1mm, font=\scriptsize\sffamily,
  st/.style={draw, rounded corners=1.5pt, fill=ln-box, align=center,
    minimum width=14.6mm, minimum height=9mm, inner sep=0.5pt},
  num/.style={font=\tiny\sffamily, text=ln-muted, anchor=south},
  net/.style={draw=none, fill=ln-accent!60},
  dsk/.style={draw=none, fill=ln-caution!60},
  >={Stealth[length=1.2mm]}]
  \node[st] (s1) at (7.3,0)   {\iflangja{接続の\\受け付け}{accept\\connection}};
  \node[st] (s2) at (24.2,0)  {\iflangja{要求の\\読み込み}{read\\request}};
  \node[st] (s3) at (41.1,0)  {\iflangja{要求の\\解析}{parse\\request}};
  \node[st] (s4) at (58.0,0)  {\iflangja{ファイル\\の検索}{find\\file}};
  \node[st] (s5) at (74.9,0)  {\iflangja{ヘッダ\\の計算}{compute\\header}};
  \node[st] (s6) at (91.8,0)  {\iflangja{ヘッダ\\の送信}{send\\header}};
  \node[st] (s7) at (108.7,0) {\iflangja{ファイル\\読込・送信}{read file,\\send data}};
  \foreach \i [evaluate=\i as \j using int(\i+1)] in {1,...,6} {
    \draw[->] (s\i.east) -- (s\j.west);
  }
  \foreach \i in {1,...,7} { \node[num] at (s\i.north) {\i}; }
  % may wait for the network
  \foreach \x in {7.3,24.2,91.8,108.7} {
    \fill[net] (\x-7.3,-6.2) rectangle (\x+7.3,-5.0);
  }
  % may wait for the disk
  \foreach \x in {58.0,108.7} {
    \fill[dsk] (\x-7.3,-8.2) rectangle (\x+7.3,-7.0);
  }
  % legend
  \fill[net] (22,-12.6) rectangle (27,-11.4);
  \node[anchor=west] at (28,-12) {\iflangja{ネットワークを待つ可能性}{may wait for the network}};
  \fill[dsk] (68,-12.6) rectangle (73,-11.4);
  \node[anchor=west] at (74,-12) {\iflangja{ディスクを待つ可能性}{may wait for the disk}};
\end{tikzpicture}

  \caption{静的な内容を提供する単純なウェブサーバの処理ステップ．ステップ
    の下の帯は，そのステップがどの資源を待つ可能性があるかを示す．}
  \label{fig:l06-web-steps}
\end{figure}

これらのステップを一度に1つの要求について実行するサーバは，1つの
クライアントにしか応じられず，自らがブロックしている間でさえ他のすべての
クライアントを待たせる．最も単純な対策は，そのコピーを複数実行することで
ある．

\begin{definition}[マルチプロセスモデル]
  \term[まるちぷろせすもでる]{マルチプロセスモデル}[multi-process model]%
  （MP）では，サーバは複数のプロセスからなり，各プロセスは，一度に1つの
  要求についてすべてのステップを実行する単一スレッドのサーバを動かす
  （\cref{fig:l06-mp-mt}a）．
\end{definition}

マルチプロセスモデルは，各プロセスが普通の逐次プログラムであるため，
プログラムしやすい．その代償はアドレス空間が分かれていることから生じる．
多数のプロセスは多くのメモリを必要とし，その分ファイルをキャッシュする
メモリが減る．プロセス間の切り替えは高価である．ファイルデータや応答
ヘッダのキャッシュのように，プロセスが共有すべき状態は，プロセス間通信か
明示的な共有メモリを必要とするため，維持するのが難しく高価である．最後に，
すべてのプロセスが同じポートで接続を受け付けなければならず，その準備には
いくらかの注意を要する．典型的には，親プロセスが待ち受けソケットを作って
からサーバプロセスを fork し，各サーバプロセスはそのソケットを継承して，
いずれもそれに対して \texttt{accept} を呼ぶ．

\begin{definition}[マルチスレッドモデル]
  \term[まるちすれっどもでる]{マルチスレッドモデル}[multithreaded model]%
  （MT）では，サーバは1つのプロセスであり，1つのアドレス空間を共有する
  そのスレッドが，それぞれ並行に要求を処理する（\cref{fig:l06-mp-mt}b）．
\end{definition}

スレッドの構成方法はいくつかある．すべてのスレッドが自分で接続を受け付け，
すべてのステップを実行してもよい．あるいは第3章のボス・ワーカーパターンの
ように，ボススレッドが接続を受け付けてワーカーに渡してもよい．あるいは各
ステップをパイプラインのステージとしてもよい．どの構成でもスレッドは
アドレス空間を共有するので，キャッシュのような共有状態に直接アクセスでき，
メモリ使用量は小さく，スレッド間のコンテキストスイッチは安価である．その
代わり実装は単純ではない．共有状態は同期を必要とし，それに伴うオーバー
ヘッドと誤りの危険を伴う．このモデルはまた，オペレーティングシステムが
カーネルレベルスレッドに対応していることに依存する．これらの
アーキテクチャが最初に比較された当時は普遍的ではなかったが，今日では
問題にならない．

\begin{figure}[htbp]
  \centering
  % Multi-process and multithreaded web servers.  Each row of seven cells is
% the sequence of processing steps of one request.
% Usage: % Multi-process and multithreaded web servers.  Each row of seven cells is
% the sequence of processing steps of one request.
% Usage: % Multi-process and multithreaded web servers.  Each row of seven cells is
% the sequence of processing steps of one request.
% Usage: \input{figures/l06-mp-mt}   (width ~116 mm)
\begin{tikzpicture}[x=1mm, y=1mm, font=\footnotesize\sffamily,
  proc/.style={draw, rounded corners=2pt, fill=white},
  cell/.style={draw, fill=ln-box, minimum width=3.6mm, minimum height=4mm,
    inner sep=0pt},
  cache/.style={draw, fill=ln-accent!15, align=center, inner sep=0.5pt,
    font=\scriptsize\sffamily},
  lab/.style={font=\scriptsize\sffamily},
  ttl/.style={font=\footnotesize\sffamily\bfseries, anchor=west},
  >={Stealth[length=1.4mm]}]
  % (a) multi-process.  The Japanese label of the cache box is wider than the
  % English one, so the cells start further left and the box sits further in.
  \iflangja{\def\lnMPcell{12.9}\def\lnMPcache{47.9}}%
           {\def\lnMPcell{13.8}\def\lnMPcache{48}}%
  \node[ttl] at (0,31) {\iflangja{(a) マルチプロセス}{(a) multi-process}};
  \foreach \y/\n in {22/1,12/2,2/3} {
    \draw[proc] (0,\y-4) rectangle (54,\y+4);
    \node[lab] at (5.5,\y) {P\n};
    \foreach \i in {0,...,6} { \node[cell] at (\lnMPcell+4.4*\i,\y) {}; }
    \node[cache, minimum width=11mm, minimum height=6mm] at (\lnMPcache,\y)
      {\iflangja{キャッシュ}{cache}};
  }
  \node[lab, text=ln-muted, anchor=north] at (27,-2.8)
    {\iflangja{プロセスごとに別のアドレス空間}{a separate address space per process}};
  % (b) multithreaded.  Same adjustment for the wider Japanese shared-state box.
  \iflangja{\def\lnMTsh{107.6}\def\lnMTlab{65.5}\def\lnMTcell{70.1}\def\lnMTarr{98.3}}%
           {\def\lnMTsh{109.5}\def\lnMTlab{66.5}\def\lnMTcell{73.8}\def\lnMTarr{102}}%
  \node[ttl] at (62,31) {\iflangja{(b) マルチスレッド}{(b) multithreaded}};
  \draw[proc] (62,-2) rectangle (116,26);
  \node[cache, minimum width=11mm, minimum height=17mm] (sh) at (\lnMTsh,12)
    {\iflangja{共有\\状態\\（キャッシュ）}{shared\\state\\(cache)}};
  \foreach \y/\n in {19/1,12/2,5/3} {
    \node[lab] at (\lnMTlab,\y) {T\n};
    \foreach \i in {0,...,6} { \node[cell] at (\lnMTcell+4.4*\i,\y) {}; }
    \draw[<->] (\lnMTarr,\y) -- (sh.west |- 0,\y);
  }
  \node[lab, text=ln-muted, anchor=north] at (89,-2.8)
    {\iflangja{1つのアドレス空間を共有}{one shared address space}};
\end{tikzpicture}
   (width ~116 mm)
\begin{tikzpicture}[x=1mm, y=1mm, font=\footnotesize\sffamily,
  proc/.style={draw, rounded corners=2pt, fill=white},
  cell/.style={draw, fill=ln-box, minimum width=3.6mm, minimum height=4mm,
    inner sep=0pt},
  cache/.style={draw, fill=ln-accent!15, align=center, inner sep=0.5pt,
    font=\scriptsize\sffamily},
  lab/.style={font=\scriptsize\sffamily},
  ttl/.style={font=\footnotesize\sffamily\bfseries, anchor=west},
  >={Stealth[length=1.4mm]}]
  % (a) multi-process.  The Japanese label of the cache box is wider than the
  % English one, so the cells start further left and the box sits further in.
  \iflangja{\def\lnMPcell{12.9}\def\lnMPcache{47.9}}%
           {\def\lnMPcell{13.8}\def\lnMPcache{48}}%
  \node[ttl] at (0,31) {\iflangja{(a) マルチプロセス}{(a) multi-process}};
  \foreach \y/\n in {22/1,12/2,2/3} {
    \draw[proc] (0,\y-4) rectangle (54,\y+4);
    \node[lab] at (5.5,\y) {P\n};
    \foreach \i in {0,...,6} { \node[cell] at (\lnMPcell+4.4*\i,\y) {}; }
    \node[cache, minimum width=11mm, minimum height=6mm] at (\lnMPcache,\y)
      {\iflangja{キャッシュ}{cache}};
  }
  \node[lab, text=ln-muted, anchor=north] at (27,-2.8)
    {\iflangja{プロセスごとに別のアドレス空間}{a separate address space per process}};
  % (b) multithreaded.  Same adjustment for the wider Japanese shared-state box.
  \iflangja{\def\lnMTsh{107.6}\def\lnMTlab{65.5}\def\lnMTcell{70.1}\def\lnMTarr{98.3}}%
           {\def\lnMTsh{109.5}\def\lnMTlab{66.5}\def\lnMTcell{73.8}\def\lnMTarr{102}}%
  \node[ttl] at (62,31) {\iflangja{(b) マルチスレッド}{(b) multithreaded}};
  \draw[proc] (62,-2) rectangle (116,26);
  \node[cache, minimum width=11mm, minimum height=17mm] (sh) at (\lnMTsh,12)
    {\iflangja{共有\\状態\\（キャッシュ）}{shared\\state\\(cache)}};
  \foreach \y/\n in {19/1,12/2,5/3} {
    \node[lab] at (\lnMTlab,\y) {T\n};
    \foreach \i in {0,...,6} { \node[cell] at (\lnMTcell+4.4*\i,\y) {}; }
    \draw[<->] (\lnMTarr,\y) -- (sh.west |- 0,\y);
  }
  \node[lab, text=ln-muted, anchor=north] at (89,-2.8)
    {\iflangja{1つのアドレス空間を共有}{one shared address space}};
\end{tikzpicture}
   (width ~116 mm)
\begin{tikzpicture}[x=1mm, y=1mm, font=\footnotesize\sffamily,
  proc/.style={draw, rounded corners=2pt, fill=white},
  cell/.style={draw, fill=ln-box, minimum width=3.6mm, minimum height=4mm,
    inner sep=0pt},
  cache/.style={draw, fill=ln-accent!15, align=center, inner sep=0.5pt,
    font=\scriptsize\sffamily},
  lab/.style={font=\scriptsize\sffamily},
  ttl/.style={font=\footnotesize\sffamily\bfseries, anchor=west},
  >={Stealth[length=1.4mm]}]
  % (a) multi-process.  The Japanese label of the cache box is wider than the
  % English one, so the cells start further left and the box sits further in.
  \iflangja{\def\lnMPcell{12.9}\def\lnMPcache{47.9}}%
           {\def\lnMPcell{13.8}\def\lnMPcache{48}}%
  \node[ttl] at (0,31) {\iflangja{(a) マルチプロセス}{(a) multi-process}};
  \foreach \y/\n in {22/1,12/2,2/3} {
    \draw[proc] (0,\y-4) rectangle (54,\y+4);
    \node[lab] at (5.5,\y) {P\n};
    \foreach \i in {0,...,6} { \node[cell] at (\lnMPcell+4.4*\i,\y) {}; }
    \node[cache, minimum width=11mm, minimum height=6mm] at (\lnMPcache,\y)
      {\iflangja{キャッシュ}{cache}};
  }
  \node[lab, text=ln-muted, anchor=north] at (27,-2.8)
    {\iflangja{プロセスごとに別のアドレス空間}{a separate address space per process}};
  % (b) multithreaded.  Same adjustment for the wider Japanese shared-state box.
  \iflangja{\def\lnMTsh{107.6}\def\lnMTlab{65.5}\def\lnMTcell{70.1}\def\lnMTarr{98.3}}%
           {\def\lnMTsh{109.5}\def\lnMTlab{66.5}\def\lnMTcell{73.8}\def\lnMTarr{102}}%
  \node[ttl] at (62,31) {\iflangja{(b) マルチスレッド}{(b) multithreaded}};
  \draw[proc] (62,-2) rectangle (116,26);
  \node[cache, minimum width=11mm, minimum height=17mm] (sh) at (\lnMTsh,12)
    {\iflangja{共有\\状態\\（キャッシュ）}{shared\\state\\(cache)}};
  \foreach \y/\n in {19/1,12/2,5/3} {
    \node[lab] at (\lnMTlab,\y) {T\n};
    \foreach \i in {0,...,6} { \node[cell] at (\lnMTcell+4.4*\i,\y) {}; }
    \draw[<->] (\lnMTarr,\y) -- (sh.west |- 0,\y);
  }
  \node[lab, text=ln-muted, anchor=north] at (89,-2.8)
    {\iflangja{1つのアドレス空間を共有}{one shared address space}};
\end{tikzpicture}

  \caption{(a) マルチプロセスモデル．各プロセスが1つの要求のすべての処理
    ステップ（マス）を実行し，キャッシュのような状態を自分の複製として
    保持する．(b) マルチスレッドモデル．1つのアドレス空間の中のスレッドが
    各ステップを実行し，状態を共有する．}
  \label{fig:l06-mp-mt}
\end{figure}

\needspace{6\baselineskip}
\section{イベント駆動モデル}
\label{sec:l06-event}

第3のアーキテクチャは，複数の実行コンテキストを用いずに並行性を実現する．
\source{P2L5，イベント駆動モデル；OSTEP ch.~33；Kerrisk \cite{kerrisk2010}
ch.~63．}

\begin{definition}[イベント駆動モデル]
  1つのアドレス空間の中で単一の制御の流れを持つ単一のプロセスとして動作し，
  イベントに応じて仕事を行うアプリケーションは，
  \term[いべんとくどうもでる]{イベント駆動モデル}[event-driven model]に
  従う．
\end{definition}

このようなアプリケーションの中心にあるのが
\term[いべんとでぃすぱっちゃ]{イベントディスパッチャ}[event dispatcher]で
ある．これは，イベントを繰り返し待ち，イベントに応じて，それに登録された
ハンドラを1つ以上呼び出すループである（\cref{fig:l06-event-driven}）．
ウェブサーバの場合，イベントには接続要求の到着，接続上での要求データの
到着，送信の完了（その接続がさらにデータを受け取れることを意味する），
そしてディスク読み込みの完了が含まれる．Pai らは，このように構築された
サーバを\term[たんいつぷろせすいべんとくどうもでる]{単一プロセス%
イベント駆動モデル}[single-process event-driven model]%
\index{SPED|see{単一プロセスイベント駆動モデル}}（SPED）の一例と呼ぶ．

\begin{figure}[htbp]
  \centering
  % Event-driven web server: one dispatcher loop invokes the handler for each
% event and keeps the state of every request explicitly.
% Usage: % Event-driven web server: one dispatcher loop invokes the handler for each
% event and keeps the state of every request explicitly.
% Usage: % Event-driven web server: one dispatcher loop invokes the handler for each
% event and keeps the state of every request explicitly.
% Usage: \input{figures/l06-event-driven}   (width ~116 mm)
\begin{tikzpicture}[x=1mm, y=1mm, font=\footnotesize\sffamily,
  box/.style={draw, rounded corners=2pt, fill=ln-box, align=center, inner sep=1pt},
  hd/.style={box, minimum width=36mm, minimum height=6mm, font=\scriptsize\sffamily},
  lab/.style={font=\scriptsize\sffamily, text=ln-muted, align=center},
  >={Stealth[length=1.6mm]}]
  % event sources
  \node[box, fill=white, minimum width=17mm, minimum height=9mm, font=\scriptsize\sffamily]
    (net) at (8.5,12) {\iflangja{ネットワーク\\（ソケット）}{network\\(sockets)}};
  \node[box, fill=white, minimum width=17mm, minimum height=9mm, font=\scriptsize\sffamily]
    (dsk) at (8.5,-4) {\iflangja{ディスク\\（ファイル）}{disk\\(files)}};
  % process
  \draw[ln-rule, rounded corners=3pt] (22,-20) rectangle (116,24);
  \node[lab, anchor=north west] at (23,23)
    {\iflangja{1 プロセス・1 スレッド・1 アドレス空間}{one process, one thread, one address space}};
  \node[box, fill=ln-accent!15, minimum width=30mm, minimum height=14mm]
    (disp) at (43,6) {\iflangja{イベントディスパッチャ}{event dispatcher}\\[0.5mm]
      \scriptsize\iflangja{イベントを待つループ}{loop: wait for events}\\
      \scriptsize(\texttt{epoll}, \texttt{select}, \texttt{poll})};
  \draw[->] (net.east) -- node[lab, above, sloped] {FD} (disp.west |- 0,9);
  \draw[->] (dsk.east) -- node[lab, below, sloped] {FD} (disp.west |- 0,3);
  \node[lab, text=ln-muted, anchor=north, align=center] at (10,-9.5)
    {\iflangja{非同期I/O\\またはヘルパ経由}{asynchronous I/O\\or a helper}};
  % handlers
  \node[hd] (h1) at (96,16) {\iflangja{接続の受け付け}{accept connection}};
  \node[hd] (h2) at (96,8)  {\iflangja{要求の読み込みと解析}{read and parse request}};
  \node[hd] (h3) at (96,0)  {\iflangja{ヘッダの送信}{send header}};
  \node[hd] (h4) at (96,-8) {\iflangja{ファイル読込・データ送信}{read file, send data}};
  \foreach \h in {h1,h2,h3,h4} { \draw[->] (disp.east) -- (\h.west); }
  \node[lab, anchor=north] at (95,-11.8)
    {\iflangja{ハンドラは最後まで実行して戻る}{handlers run to completion and return}};
  % per-request state
  \node[box, fill=white, minimum width=30mm, align=left, font=\scriptsize\sffamily,
    anchor=north] (st) at (43,-4.5)
    {C1: \iflangja{要求を待つ}{waiting for request}\\
     C2: \iflangja{送信を待つ}{waiting to send}\\
     C3: \iflangja{ファイルを待つ}{waiting for file data}};
  \node[lab, anchor=north] at (st.south) {\iflangja{要求ごとの状態}{per-request state}};
\end{tikzpicture}
   (width ~116 mm)
\begin{tikzpicture}[x=1mm, y=1mm, font=\footnotesize\sffamily,
  box/.style={draw, rounded corners=2pt, fill=ln-box, align=center, inner sep=1pt},
  hd/.style={box, minimum width=36mm, minimum height=6mm, font=\scriptsize\sffamily},
  lab/.style={font=\scriptsize\sffamily, text=ln-muted, align=center},
  >={Stealth[length=1.6mm]}]
  % event sources
  \node[box, fill=white, minimum width=17mm, minimum height=9mm, font=\scriptsize\sffamily]
    (net) at (8.5,12) {\iflangja{ネットワーク\\（ソケット）}{network\\(sockets)}};
  \node[box, fill=white, minimum width=17mm, minimum height=9mm, font=\scriptsize\sffamily]
    (dsk) at (8.5,-4) {\iflangja{ディスク\\（ファイル）}{disk\\(files)}};
  % process
  \draw[ln-rule, rounded corners=3pt] (22,-20) rectangle (116,24);
  \node[lab, anchor=north west] at (23,23)
    {\iflangja{1 プロセス・1 スレッド・1 アドレス空間}{one process, one thread, one address space}};
  \node[box, fill=ln-accent!15, minimum width=30mm, minimum height=14mm]
    (disp) at (43,6) {\iflangja{イベントディスパッチャ}{event dispatcher}\\[0.5mm]
      \scriptsize\iflangja{イベントを待つループ}{loop: wait for events}\\
      \scriptsize(\texttt{epoll}, \texttt{select}, \texttt{poll})};
  \draw[->] (net.east) -- node[lab, above, sloped] {FD} (disp.west |- 0,9);
  \draw[->] (dsk.east) -- node[lab, below, sloped] {FD} (disp.west |- 0,3);
  \node[lab, text=ln-muted, anchor=north, align=center] at (10,-9.5)
    {\iflangja{非同期I/O\\またはヘルパ経由}{asynchronous I/O\\or a helper}};
  % handlers
  \node[hd] (h1) at (96,16) {\iflangja{接続の受け付け}{accept connection}};
  \node[hd] (h2) at (96,8)  {\iflangja{要求の読み込みと解析}{read and parse request}};
  \node[hd] (h3) at (96,0)  {\iflangja{ヘッダの送信}{send header}};
  \node[hd] (h4) at (96,-8) {\iflangja{ファイル読込・データ送信}{read file, send data}};
  \foreach \h in {h1,h2,h3,h4} { \draw[->] (disp.east) -- (\h.west); }
  \node[lab, anchor=north] at (95,-11.8)
    {\iflangja{ハンドラは最後まで実行して戻る}{handlers run to completion and return}};
  % per-request state
  \node[box, fill=white, minimum width=30mm, align=left, font=\scriptsize\sffamily,
    anchor=north] (st) at (43,-4.5)
    {C1: \iflangja{要求を待つ}{waiting for request}\\
     C2: \iflangja{送信を待つ}{waiting to send}\\
     C3: \iflangja{ファイルを待つ}{waiting for file data}};
  \node[lab, anchor=north] at (st.south) {\iflangja{要求ごとの状態}{per-request state}};
\end{tikzpicture}
   (width ~116 mm)
\begin{tikzpicture}[x=1mm, y=1mm, font=\footnotesize\sffamily,
  box/.style={draw, rounded corners=2pt, fill=ln-box, align=center, inner sep=1pt},
  hd/.style={box, minimum width=36mm, minimum height=6mm, font=\scriptsize\sffamily},
  lab/.style={font=\scriptsize\sffamily, text=ln-muted, align=center},
  >={Stealth[length=1.6mm]}]
  % event sources
  \node[box, fill=white, minimum width=17mm, minimum height=9mm, font=\scriptsize\sffamily]
    (net) at (8.5,12) {\iflangja{ネットワーク\\（ソケット）}{network\\(sockets)}};
  \node[box, fill=white, minimum width=17mm, minimum height=9mm, font=\scriptsize\sffamily]
    (dsk) at (8.5,-4) {\iflangja{ディスク\\（ファイル）}{disk\\(files)}};
  % process
  \draw[ln-rule, rounded corners=3pt] (22,-20) rectangle (116,24);
  \node[lab, anchor=north west] at (23,23)
    {\iflangja{1 プロセス・1 スレッド・1 アドレス空間}{one process, one thread, one address space}};
  \node[box, fill=ln-accent!15, minimum width=30mm, minimum height=14mm]
    (disp) at (43,6) {\iflangja{イベントディスパッチャ}{event dispatcher}\\[0.5mm]
      \scriptsize\iflangja{イベントを待つループ}{loop: wait for events}\\
      \scriptsize(\texttt{epoll}, \texttt{select}, \texttt{poll})};
  \draw[->] (net.east) -- node[lab, above, sloped] {FD} (disp.west |- 0,9);
  \draw[->] (dsk.east) -- node[lab, below, sloped] {FD} (disp.west |- 0,3);
  \node[lab, text=ln-muted, anchor=north, align=center] at (10,-9.5)
    {\iflangja{非同期I/O\\またはヘルパ経由}{asynchronous I/O\\or a helper}};
  % handlers
  \node[hd] (h1) at (96,16) {\iflangja{接続の受け付け}{accept connection}};
  \node[hd] (h2) at (96,8)  {\iflangja{要求の読み込みと解析}{read and parse request}};
  \node[hd] (h3) at (96,0)  {\iflangja{ヘッダの送信}{send header}};
  \node[hd] (h4) at (96,-8) {\iflangja{ファイル読込・データ送信}{read file, send data}};
  \foreach \h in {h1,h2,h3,h4} { \draw[->] (disp.east) -- (\h.west); }
  \node[lab, anchor=north] at (95,-11.8)
    {\iflangja{ハンドラは最後まで実行して戻る}{handlers run to completion and return}};
  % per-request state
  \node[box, fill=white, minimum width=30mm, align=left, font=\scriptsize\sffamily,
    anchor=north] (st) at (43,-4.5)
    {C1: \iflangja{要求を待つ}{waiting for request}\\
     C2: \iflangja{送信を待つ}{waiting to send}\\
     C3: \iflangja{ファイルを待つ}{waiting for file data}};
  \node[lab, anchor=north] at (st.south) {\iflangja{要求ごとの状態}{per-request state}};
\end{tikzpicture}

  \caption{イベント駆動のウェブサーバ．ディスパッチャはファイル
    ディスクリプタ（FD）上のイベント---準備の整ったソケット，および非同期
    I/O やヘルパ（\cref{sec:l06-helpers}）が完了を通知するディスク読み込み
    ---を待ち，イベントごとにそのハンドラを呼び出す．処理中の各要求の状態
    は，明示的なレコードに保持される．}
  \label{fig:l06-event-driven}
\end{figure}

\term[いべんとはんどら]{イベントハンドラ}[event handler]とは，1つの
イベントに応じて実行されるコードであり，その呼び出しは単にそのコードを
呼ぶことにすぎない．ハンドラは最後まで実行される．%
\caution{最後まで実行するとは，ブロックする呼び出しを一切含まないという
ことである．\texttt{getaddrinfo} による名前解決のようにライブラリ関数に
隠れたものも，1つではなく全部の接続を止める．}%
ブロックすることになる場合，ハンドラは待たない．ブロッキング操作を開始し，要求がどこまで進んだか
を記録して，制御をディスパッチャに返す．操作の完了は後に別のイベントとして
届く．ディスパッチャとハンドラは合わせて状態機械を形成する．スレッドで
あればプログラムカウンタとスタックに暗黙のうちに保持する情報が，ここでは
要求ごとのレコードに明示的に保持される．\margin{OSTEP はこの明示的な管理を
\emph{手動スタック管理}と呼んでいる．}

\subsection{並行実行}

単一のスレッドしか持たないにもかかわらず，イベント駆動のサーバは多数の
要求を並行に処理する．処理中の各要求は，その処理のいずれかのステップに
ある．ディスパッチャは1つの要求の1ステップを実行し，それから次のイベントに
移る．そのイベントは別の要求のものかもしれず，こうして要求の処理は互いに
織り込まれる．\cref{tab:l06-trace}は3つの接続をたどったものである．例えば
ステップ 6 の後には，1つのスレッドが3つの要求を処理中で，それぞれ異なる
ステップにある．この追跡では，ファイル読み込みの完了がイベントとして届く
ことを仮定しており，そのためには非同期I/O か\cref{sec:l06-helpers}のヘルパ
が必要である．ファイルの検索もディスクを待つ場合がある．

\begin{table}[htbp]
  \centering
  \caption{3つの接続を持つイベント駆動サーバの追跡．状態は，R が要求待ち，
    F がファイルデータ待ち，S が送信可能になるのを待つ状態，--- が未接続を
    表す．}
  \label{tab:l06-trace}
  \small
  \begin{tabular}{@{}c>{\raggedright\arraybackslash}p{0.25\linewidth}>{\raggedright\arraybackslash}p{0.33\linewidth}ccc@{}}
    \toprule
    ステップ & イベント & ハンドラの動作 & C1 & C2 & C3 \\
    \midrule
    1 & 接続要求 & C1 を受け付ける & R & --- & --- \\
    2 & 接続要求 & C2 を受け付ける & R & R & --- \\
    3 & C1 上の要求 & 解析，ファイル検索，読み込み開始 & F & R & --- \\
    4 & 接続要求 & C3 を受け付ける & F & R & R \\
    5 & C2 上の要求 & 解析，ファイル検索，読み込み開始 & F & F & R \\
    6 & C1 のファイルデータ & ヘッダとデータを送信 & S & F & R \\
    7 & C3 上の要求 & 解析，ファイル検索，読み込み開始 & S & F & F \\
    8 & C1 がデータを受け取れる & 残りを送信して閉じる & 完了 & F & F \\
    9 & C2 のファイルデータ & ヘッダとデータを送信 & 完了 & S & F \\
    \bottomrule
  \end{tabular}
\end{table}

\subsection{このモデルが機能する理由}

第3章では，単一の CPU 上で，あるスレッドが I/O を待つ間に別のスレッドへ
切り替えることが，アイドル時間が2回のコンテキストスイッチのコストを上回る
とき，すなわち $t_{\mathrm{idle}} > 2\, t_{\mathrm{ctx}}$ のときに引き合う
ことを示した．アイドル時間がなければ，コンテキストスイッチは純粋な
オーバーヘッドである．要求を処理できたはずの CPU サイクルを消費するから
である．%
\margin{第2章にその費用がある．1回の切り替えの直接の仕事はおよそ
\SI{3.8}{\micro\second}，キャッシュを詰め直す分まで含めると
\SI{1000}{\micro\second} を超える \cite{li2007}．}%
イベント駆動モデルは，このコストなしに利点を得る．要求が待たなけれ
ばならなくなるまでその要求を処理し，それから別の要求を処理するのであって，
コンテキストスイッチは介在しない．複数の CPU を持つマシンでは，このモデルを
CPU ごとに1つ適用する．すなわち，それぞれが独自のディスパッチャを持つ複数
のイベント駆動プロセスが，到着する接続を分担する．

\subsection{実装}

ウェブサーバのイベントの発生源には，ファイルディスクリプタを通じて
アクセスする．
ネットワークにはソケット，ディスクにはファイル，そして以下で導入するヘルパ
にはパイプである．ソケットやパイプの場合，イベントとはそのディスクリプタが
入力または出力の準備を整えたことである．通常のファイルの場合，準備が
整っているということは，読み込みがディスクを待たなければならないかどうかに
ついて何も語らない．この制限には\cref{sec:l06-helpers}で立ち返る．どの
ディスクリプタの準備が整っているかを知るために，ディスパッチャは3つの
システムコールのいずれかを使う．\texttt{select} と \texttt{poll} は，呼び
出しのたびに関心のあるディスクリプタの集合全体を受け取る．カーネルはその
すべてを調べ，アプリケーションは結果を走査するが，準備が整っているのは
たいていごく少数である．したがってそのコストは，開いている接続の数とともに
増える．%
\caution{\texttt{select} には上限もある．Linux では \texttt{fd\_set} が
保持できるのは \texttt{FD\_SETSIZE} = \num{1024} 個のディスクリプタである．
\texttt{poll} と \texttt{epoll} にこの制限はない．}%
Linux で利用できる \texttt{epoll} は，このオーバーヘッドの大半を
取り除く．ディスクリプタは一度登録すればよく，各呼び出しは準備の整った
ものだけを返す．%
\sidenote{1台のマシンで1万の接続を扱うことは，Kegel が 1999 年にその手法を
まとめたことにちなんで \emph{C10K 問題}と呼ばれた．答えがこれらの
スケーラブルなインタフェースであり，FreeBSD~4.1 の \texttt{kqueue}
（2000 年）と Linux~2.5.44 の \texttt{epoll}（2002 年，一般に使えるように
なったのは Linux~2.6）である．}

\begin{example}
  \texttt{epoll} に基づくディスパッチャの骨組み．各要求の状態は
  \texttt{struct conn} に保持する．
  \begin{lstlisting}[language=C, numbers=none]
struct conn {                   /* 要求ごとに1つ */
    int fd;
    void (*handler)(struct conn *);   /* 次のステップ */
    /* バッファ，ファイル，オフセットなど */
};

int ep = epoll_create1(0);
/* ハンドラが epoll_ctl でディスクリプタを登録 */
for (;;) {                                /* ディスパッチャ */
    struct epoll_event ev[64];
    int n = epoll_wait(ep, ev, 64, -1);   /* 待つ */
    for (int i = 0; i < n; i++) {
        struct conn *c = ev[i].data.ptr;
        c->handler(c);           /* 最後まで実行される */
    }
}
\end{lstlisting}
  待ち受けソケットのハンドラは接続を受け付け，要求を読むハンドラを持つ
  \texttt{struct conn} を確保し，新しいディスクリプタを登録する．どの
  ハンドラも1つのステップを実行し，\texttt{handler} を次のステップに
  設定する．%
  \margin{\texttt{EPOLLET} を使うと \texttt{epoll} は状態が変わるたびに
  1度だけディスクリプタを報告するので，ハンドラは \texttt{EAGAIN} が返る
  まで読み書きしなければならない．}
\end{example}

このモデルの利点は，単一のアドレス空間と単一の制御の流れから生じる．処理中
の1つの要求が要するのはその状態レコードだけであり，スタックを伴うスレッド
ではないので，必要なメモリは小さい．実行コンテキスト間のコンテキスト
スイッチは存在しない．そしてハンドラが並行に実行されることはないので，同期
は不要である．

\begin{keypoint}
  イベント駆動のサーバは，多数の要求の処理を1つのスレッドの中で互いに
  織り込む．ディスパッチャがファイルディスクリプタ上のイベントを待ち，
  それぞれについてハンドラを実行する．どのハンドラもブロックしない限り，
  メモリ，コンテキストスイッチ，同期を節約できる．
\end{keypoint}

\needspace{6\baselineskip}
\section{ブロッキング操作とヘルパ}
\label{sec:l06-helpers}

イベント駆動モデルには1つ重大な問題がある．ブロックするハンドラがプロセス
全体をブロックしてしまうことである．\source{P2L5，非同期I/O，ヘルパ；
OSTEP ch.~33；Pai ら．}単一のスレッドがディスクを待つ間，ディスパッチャは
実行されず，他のどの要求も前に進まない．データを送る準備が整っている要求
さえも進まない．準備完了の通知はディスク読み込みの役に立たない．
\texttt{select} と \texttt{poll} は，通常のファイルを読むのにディスクを
待たなければならない場合でも常に準備が整っていると報告し，\texttt{epoll}
は通常のファイルをそもそも受け付けない．

\begin{definition}[非同期I/O]
  \term[ひどうきIO]{非同期I/O}[asynchronous I/O]では，I/O 操作を開始する
  システムコールがその操作の完了前に戻り，呼び出したスレッドは実行を続ける．
  カーネルはその呼び出しから，操作を遂行するのに必要な情報をすべて得て，
  結果をどこへ届ければよいかを知るか，あるいは後で結果をどこから取り出せば
  よいかを呼び出し側に伝える．
\end{definition}

非同期I/O はカーネルからの支援を必要とする．例えば，操作のブロックする部分
を実行するカーネルレベルスレッドである．あるいはデバイスからの支援を必要と
する．例えば，DMA（第11章）でデータを転送し，完了を割り込みで知らせる
デバイスである．非同期I/O はイベント駆動モデルによく適合する．非同期操作の
完了は，ディスパッチャにとってもう1つのイベントにすぎないからである．
\margin{POSIX は \texttt{aio\_read} と
\texttt{aio\_write} を定めており，glibc はこれらをスレッドで実装している．
Linux でバッファ付きのファイル I/O まで扱えるのは \texttt{io\_uring}
（2019 年）である．}

しかし非同期の呼び出しは常に利用できるわけではなく，存在する場合でも
あらゆる種類のデバイスについて対応しているとは限らない．そのときイベント
駆動のサーバは，ブロッキング操作を委譲する．

\begin{definition}[ヘルパ]
  \term[へるぱ]{ヘルパ}[helper]とは，イベント駆動のサーバが，自らの
  ディスパッチャに代わってブロッキング操作を実行させるためだけに用いる
  スレッドまたはプロセスである．
\end{definition}

ディスパッチャは操作をパイプでヘルパに渡し，ヘルパは完了を第2のパイプで
報告する．パイプはデータを一方向にしか運ばないからである．ソケットの対なら
両方向を一度に担える．パイプもソケットもファイルディスクリプタなので，
ディスパッチャはこれらの報告を，他のイベントと同じように \texttt{select}
や \texttt{poll} で待つ．ブロックするのはヘルパであって，ディスパッチャは
ブロックせず，したがって他のすべての要求も止まらない．イベント駆動の
プロセスとヘルパプロセスを組み合わせたものが，
\term[ひたいしょうまるちぷろせすいべんとくどうもでる]{非対称マルチプロセス%
イベント駆動モデル}[asymmetric multi-process event-driven model]%
\index{AMPED|see{非対称マルチプロセスイベント駆動モデル}}（AMPED）である．
このモデルが\emph{非対称}であるのは，MP モデルや MT モデルの同等なプロセス
やスレッドとは違い，ヘルパが特定の操作だけを行うからである．Pai らが
ヘルパにプロセスを用いたのは，当時はすべてのオペレーティングシステムが
カーネルレベルスレッドに対応していたわけではなかったからである．プロセスの
代わりにヘルパスレッドを用いれば，サーバは
\term[ひたいしょうまるちすれっどいべんとくどうもでる]{非対称マルチスレッド%
イベント駆動モデル}[asymmetric multithreaded event-driven model]%
\index{AMTED|see{非対称マルチスレッドイベント駆動モデル}}（AMTED）に従う．

ヘルパは，基本のイベント駆動モデルの移植性の問題を解消する．サーバが
カーネルの非同期I/O 支援に依存しなくなるからである．またヘルパは，要求ごと
にワーカーを置く場合よりもメモリ使用量を小さく保つ．ヘルパが必要なのは同時
にブロックしている操作の分だけであり，接続ごとではないからである．他方で
この手法は，ある種のアプリケーションにしか適用できず，マルチプロセッサ上
ではイベントの振り分けを複雑にする．ヘルパが報告する完了も含め，各イベント
は対応する要求を扱っているディスパッチャに届かなければならない．

\begin{example}
  あるサーバが 2\,000 個の接続を開いており，そのうち任意の時点でディスクを
  待っているのは高々 25 個であるとする．マルチスレッドモデルでは接続ごとに
  スレッドが必要であり，スタックとデータ構造にスレッド当たり
  \SI{64}{\kilo\byte} を要するとすれば，スレッドは \SI{128}{\mega\byte} を
  使う．%
  \margin{これは実際に使う分である．Pthreads のスタックはアドレス空間を
  \SI{8}{\mebi\byte} 予約するが，ページを使うのは触れた分だけである
  （第4章）．}%
  AMPED のサーバは，接続ごとに例えば \SI{2}{\kilo\byte} の状態
  レコード（\SI{4}{\mega\byte}）と，プロセス固有のページ，ページテーブル，
  カーネル構造に例えば1個当たり \SI{1}{\mega\byte} を要する 25 個のヘルパ
  プロセス（\SI{25}{\mega\byte}）を保持し，合計でおよそ
  \SI{29}{\mega\byte} となる．その差をファイルデータのキャッシュに使えば，
  ディスクを待たなければならない要求の数が減る．
\end{example}

\cref{tab:l06-models}は4つのアーキテクチャをまとめたものである．

\begin{table}[htbp]
  \centering
  \caption{並行ウェブサーバのアーキテクチャ．}
  \label{tab:l06-models}
  \small
  \begin{tabular}{@{}>{\raggedright\arraybackslash}p{0.19\linewidth}>{\raggedright\arraybackslash}p{0.165\linewidth}>{\raggedright\arraybackslash}p{0.165\linewidth}>{\raggedright\arraybackslash}p{0.165\linewidth}>{\raggedright\arraybackslash}p{0.165\linewidth}@{}}
    \toprule
    & MP & MT & SPED & AMPED \\
    \midrule
    実行コンテキスト & 並行する要求ごとに1プロセス & 並行する要求ごとに
      1スレッド & 1つ & 1つとヘルパ \\
    要求当たりのメモリ & アドレス空間 & スレッドのスタックと構造 & 状態
      レコード & 状態レコード \\
    コンテキストスイッチ & プロセス間 & スレッド間 & なし & ブロッキング
      操作のときだけ \\
    同期 & プロセス間通信による & 共有状態のロック & なし & なし \\
    ブロックするディスク I/O & 1プロセスがブロック & 1スレッドがブロック &
      サーバがブロック & 1つのヘルパがブロック \\
    \bottomrule
  \end{tabular}
\end{table}

\needspace{6\baselineskip}
\section{Flash ウェブサーバ}
\label{sec:l06-flash}

Pai らは，AMPED モデルに基づくイベント駆動のウェブサーバを構築した．
\source{P2L5，Flash；Pai ら；Kerrisk ch.~49--50．}彼らのサーバ
\term{Flash}[Flash] では，ブロックしうるディスク操作をヘルパが行う．
すなわち，ディレクトリの読み込みを伴う場合のある，要求のパス名から
ファイルへの対応付けと，ファイルデータの読み込みである．ヘルパは
ディスパッチャとパイプで通信する．Flash はファイルデータにメモリマップド
ファイル（第9章）を通じてアクセスするので，ファイルのページは最初に
アクセスされたときにディスクから読まれる．

ハンドラがファイルデータを送る前に，ディスパッチャはシステムコール
\texttt{mincore} を使って，そのデータのページがメモリ上にあるかどうかを
調べる（\cref{fig:l06-flash}）．%
\tip{\texttt{mincore} はページ当たり1バイトで答える．x86-64 の
\SI{4}{\kibi\byte} ページなら，\SI{4}{\mebi\byte} の対応付けに必要な
ベクタは \num{1024} バイトである．}%
メモリ上にあれば，ハンドラは直ちにデータを送り，これはブロックしえない．メモリ上になければ，ディスパッチャはその要求
をパイプでヘルパに渡す．ヘルパはそれらのページにアクセスし，それによって
ページがメモリに読み込まれるが，その間ブロックしうる．そしてヘルパは完了を
パイプでディスパッチャに報告する．ディスパッチャはその報告をイベントとして
受け取り，ハンドラを呼び出す．ハンドラは今度はデータをメモリ上に見つける．
ウェブサーバへの要求の大半は，データがすでにメモリ上にある人気のファイルに
対するものであり，そうした要求では Flash はヘルパを使わずに済み，ヘルパ
との通信のコストも避けられる．これは大きな節約になりうる．

\begin{figure}[htbp]
  \centering
  % Flash (AMPED): the dispatcher tests with mincore whether the file data are
% in memory; if not, a helper process reads them and reports completion
% through a pipe.
% Usage: % Flash (AMPED): the dispatcher tests with mincore whether the file data are
% in memory; if not, a helper process reads them and reports completion
% through a pipe.
% Usage: % Flash (AMPED): the dispatcher tests with mincore whether the file data are
% in memory; if not, a helper process reads them and reports completion
% through a pipe.
% Usage: \input{figures/l06-flash}   (width ~116 mm)
\begin{tikzpicture}[x=1mm, y=1mm, font=\footnotesize\sffamily,
  box/.style={draw, rounded corners=2pt, fill=ln-box, align=center, inner sep=1pt},
  lab/.style={font=\scriptsize\sffamily, align=center},
  num/.style={draw, circle, fill=white, inner sep=0pt, minimum size=4mm,
    font=\scriptsize\sffamily\bfseries, text=ln-accent},
  >={Stealth[length=1.6mm]}]
  % main process
  \draw[ln-rule, rounded corners=3pt] (0,-15) rectangle (64,21);
  \node[lab, text=ln-muted, anchor=north west] at (1,20)
    {\iflangja{Flash のプロセス（ディスパッチャとハンドラ）}{Flash process (dispatcher and handlers)}};
  \node[box, fill=ln-accent!15, minimum width=20mm, minimum height=10mm]
    (disp) at (12,6) {\iflangja{ディスパッチャ}{dispatcher}};
  \node[box, fill=white, minimum width=24mm, minimum height=10mm, font=\scriptsize\sffamily]
    (chk) at (48,6) {\iflangja{ページはメモリ上か}{pages in memory?}\\(\texttt{mincore})};
  \node[box, minimum width=24mm, minimum height=8mm, font=\scriptsize\sffamily]
    (hnd) at (48,-8) {\iflangja{ハンドラ：}{handler:}\\\iflangja{データを送信}{send data}};
  \draw[->] (disp) -- node[num, above=0.6mm] {1} (chk);
  \draw[->] (chk) -- node[lab, right] {\iflangja{はい}{yes}} node[num, left=0.6mm] {2} (hnd);
  % helpers
  \foreach \y/\n in {22/1,14/2,6/3} {
    \node[box, minimum width=16mm, minimum height=6mm, font=\scriptsize\sffamily]
      (h\n) at (86,\y) {\iflangja{ヘルパ}{helper} \n};
  }
  \draw[->] (chk) -- node[lab, above, pos=0.62] {\iflangja{いいえ：\\パイプ}{no: pipe}}
    node[num, below=0.6mm, pos=0.62] {3} (h3);
  % disk
  \node[cylinder, shape border rotate=90, draw, aspect=0.25, fill=ln-box,
    minimum width=11mm, minimum height=13mm, font=\scriptsize\sffamily]
    (disk) at (109,6) {\iflangja{ディスク}{disk}};
  \draw[<->] (h3.east) -- node[lab, above] {\texttt{mmap}} (disk.west);
  % completion
  \draw[->, densely dashed, ln-caution] (h3.south) -- (86,-20) -- (12,-20)
    -- node[num, left=0.6mm, pos=0.5] {4} (disp.south);
  \node[lab, text=ln-caution, anchor=south] at (48,-19.6)
    {\iflangja{完了をパイプで通知（イベント）}{completion reported through a pipe (an event)}};
  \draw[->, densely dashed, ln-caution] (disp.south east) -- (hnd.west);
\end{tikzpicture}
   (width ~116 mm)
\begin{tikzpicture}[x=1mm, y=1mm, font=\footnotesize\sffamily,
  box/.style={draw, rounded corners=2pt, fill=ln-box, align=center, inner sep=1pt},
  lab/.style={font=\scriptsize\sffamily, align=center},
  num/.style={draw, circle, fill=white, inner sep=0pt, minimum size=4mm,
    font=\scriptsize\sffamily\bfseries, text=ln-accent},
  >={Stealth[length=1.6mm]}]
  % main process
  \draw[ln-rule, rounded corners=3pt] (0,-15) rectangle (64,21);
  \node[lab, text=ln-muted, anchor=north west] at (1,20)
    {\iflangja{Flash のプロセス（ディスパッチャとハンドラ）}{Flash process (dispatcher and handlers)}};
  \node[box, fill=ln-accent!15, minimum width=20mm, minimum height=10mm]
    (disp) at (12,6) {\iflangja{ディスパッチャ}{dispatcher}};
  \node[box, fill=white, minimum width=24mm, minimum height=10mm, font=\scriptsize\sffamily]
    (chk) at (48,6) {\iflangja{ページはメモリ上か}{pages in memory?}\\(\texttt{mincore})};
  \node[box, minimum width=24mm, minimum height=8mm, font=\scriptsize\sffamily]
    (hnd) at (48,-8) {\iflangja{ハンドラ：}{handler:}\\\iflangja{データを送信}{send data}};
  \draw[->] (disp) -- node[num, above=0.6mm] {1} (chk);
  \draw[->] (chk) -- node[lab, right] {\iflangja{はい}{yes}} node[num, left=0.6mm] {2} (hnd);
  % helpers
  \foreach \y/\n in {22/1,14/2,6/3} {
    \node[box, minimum width=16mm, minimum height=6mm, font=\scriptsize\sffamily]
      (h\n) at (86,\y) {\iflangja{ヘルパ}{helper} \n};
  }
  \draw[->] (chk) -- node[lab, above, pos=0.62] {\iflangja{いいえ：\\パイプ}{no: pipe}}
    node[num, below=0.6mm, pos=0.62] {3} (h3);
  % disk
  \node[cylinder, shape border rotate=90, draw, aspect=0.25, fill=ln-box,
    minimum width=11mm, minimum height=13mm, font=\scriptsize\sffamily]
    (disk) at (109,6) {\iflangja{ディスク}{disk}};
  \draw[<->] (h3.east) -- node[lab, above] {\texttt{mmap}} (disk.west);
  % completion
  \draw[->, densely dashed, ln-caution] (h3.south) -- (86,-20) -- (12,-20)
    -- node[num, left=0.6mm, pos=0.5] {4} (disp.south);
  \node[lab, text=ln-caution, anchor=south] at (48,-19.6)
    {\iflangja{完了をパイプで通知（イベント）}{completion reported through a pipe (an event)}};
  \draw[->, densely dashed, ln-caution] (disp.south east) -- (hnd.west);
\end{tikzpicture}
   (width ~116 mm)
\begin{tikzpicture}[x=1mm, y=1mm, font=\footnotesize\sffamily,
  box/.style={draw, rounded corners=2pt, fill=ln-box, align=center, inner sep=1pt},
  lab/.style={font=\scriptsize\sffamily, align=center},
  num/.style={draw, circle, fill=white, inner sep=0pt, minimum size=4mm,
    font=\scriptsize\sffamily\bfseries, text=ln-accent},
  >={Stealth[length=1.6mm]}]
  % main process
  \draw[ln-rule, rounded corners=3pt] (0,-15) rectangle (64,21);
  \node[lab, text=ln-muted, anchor=north west] at (1,20)
    {\iflangja{Flash のプロセス（ディスパッチャとハンドラ）}{Flash process (dispatcher and handlers)}};
  \node[box, fill=ln-accent!15, minimum width=20mm, minimum height=10mm]
    (disp) at (12,6) {\iflangja{ディスパッチャ}{dispatcher}};
  \node[box, fill=white, minimum width=24mm, minimum height=10mm, font=\scriptsize\sffamily]
    (chk) at (48,6) {\iflangja{ページはメモリ上か}{pages in memory?}\\(\texttt{mincore})};
  \node[box, minimum width=24mm, minimum height=8mm, font=\scriptsize\sffamily]
    (hnd) at (48,-8) {\iflangja{ハンドラ：}{handler:}\\\iflangja{データを送信}{send data}};
  \draw[->] (disp) -- node[num, above=0.6mm] {1} (chk);
  \draw[->] (chk) -- node[lab, right] {\iflangja{はい}{yes}} node[num, left=0.6mm] {2} (hnd);
  % helpers
  \foreach \y/\n in {22/1,14/2,6/3} {
    \node[box, minimum width=16mm, minimum height=6mm, font=\scriptsize\sffamily]
      (h\n) at (86,\y) {\iflangja{ヘルパ}{helper} \n};
  }
  \draw[->] (chk) -- node[lab, above, pos=0.62] {\iflangja{いいえ：\\パイプ}{no: pipe}}
    node[num, below=0.6mm, pos=0.62] {3} (h3);
  % disk
  \node[cylinder, shape border rotate=90, draw, aspect=0.25, fill=ln-box,
    minimum width=11mm, minimum height=13mm, font=\scriptsize\sffamily]
    (disk) at (109,6) {\iflangja{ディスク}{disk}};
  \draw[<->] (h3.east) -- node[lab, above] {\texttt{mmap}} (disk.west);
  % completion
  \draw[->, densely dashed, ln-caution] (h3.south) -- (86,-20) -- (12,-20)
    -- node[num, left=0.6mm, pos=0.5] {4} (disp.south);
  \node[lab, text=ln-caution, anchor=south] at (48,-19.6)
    {\iflangja{完了をパイプで通知（イベント）}{completion reported through a pipe (an event)}};
  \draw[->, densely dashed, ln-caution] (disp.south east) -- (hnd.west);
\end{tikzpicture}

  \caption{Flash におけるファイルデータの提供．(1)~ディスパッチャが
    \texttt{mincore} でデータがメモリ上にあるかを調べる．(2)~あれば，
    ハンドラが直接それを送る．(3)~なければ，ヘルパにパイプでページへの
    アクセスを依頼し，ヘルパがディスクから読み込む．(4)~その完了報告が
    イベントとなり，その後ハンドラがデータを送る．}
  \label{fig:l06-flash}
\end{figure}

\begin{example}
  \texttt{sys/mman.h}，\texttt{unistd.h}，\texttt{stdbool.h} を用いた C で
  の判定．ファイルデータはページ境界に整列したアドレスに対応付けられている
  ものとする．
  \begin{lstlisting}[language=C, numbers=none]
static bool in_memory(void *addr, size_t len)
{
    size_t pg = (size_t)sysconf(_SC_PAGESIZE);
    size_t n = (len + pg - 1) / pg;
    unsigned char vec[n];            /* ページ当たり1バイト */
    if (mincore(addr, len, vec) != 0)
        return false;
    for (size_t i = 0; i < n; i++)
        if (!(vec[i] & 1))           /* ページが非常駐 */
            return false;
    return true;
}

/* ディスパッチャ，ファイルデータを送る前 */
if (in_memory(c->map, c->len))
    send_data(c);                    /* ブロックしえない */
else
    ask_helper(c);                   /* パイプ経由 */
\end{lstlisting}
\end{example}

\subsection{追加の最適化}

Flash はこのアーキテクチャに，要求当たりの作業量を減らす最適化を組み合わせ
ている．
\begin{description}
  \item[アプリケーションレベルのキャッシュ] Flash はデータと計算結果の
    両方をキャッシュする．要求のパス名からファイルへの対応付け，ファイル
    だけで決まる応答ヘッダ，そして対応付けられたファイル自体である．パス名
    のキャッシュにヒットすれば，ヘルパへの往復の依頼も省ける．
  \item[ギャザ I/O] 応答ヘッダとファイルデータは，1回のベクタ I/O 操作
    \texttt{writev} で送られる．これは複数のメモリ領域を一度に渡すので
    システムコールを1回節約する．スキャッタギャザ DMA を備えたネットワーク
    インタフェースは，連続していないバッファを1つにまとめることなく転送
    できる．%
    \margin{Linux はさらに進んでいる．\texttt{sendfile}（2.2，1999 年）は
    ファイルデータをユーザ空間を経由せずカーネル内でソケットへ複写する．}
  \item[整列] カーネルは \texttt{writev} に渡された領域を自身の連続した
    バッファへコピーし，領域がマシンのワードに整列していないとそのコピーは
    遅くなる．そこで Flash は，可変長のヘッダフィールドを伸ばして応答
    ヘッダの長さを常に 32 バイトの倍数に揃え，ヘッダに続くファイルデータが
    整列した位置でコピーされるようにする．\caution{この最適化は通常
    \emph{DMA のための整列}と呼ばれるが，カーネルバッファへのコピーを避ける
    のではなく，そのコピーを安価にするものである．}
\end{description}
これらの最適化は 1999 年には新しかったが，今日のウェブサーバでは一般的で
ある．これらはサーバのアーキテクチャとは独立であり，この事実は
\cref{sec:l06-methodology}の評価において重要になる．

\needspace{6\baselineskip}
\section{Apache ウェブサーバ}
\label{sec:l06-apache}

オープンソースの \term{Apache}[Apache] HTTP サーバは，実運用のサーバが
これらのモデルをどう組み合わせるかを示している．\source{P2L5，Apache
ウェブサーバ．}その\emph{コア}がサーバの基本的な骨組みを提供し，接続を
受け付けて実行コンテキストを管理する---Apache~2 では，この節の最後の注で
述べるように，専用のモジュールを通じてそれを行う．要求を処理する機能は
\emph{モジュール}として実装され，各モジュールは要求処理の段階である
\emph{フック}にハンドラを登録する（\cref{fig:l06-apache}a）．要求が
進むにつれて，コアは各段階に登録されたハンドラを呼び出す．この制御の
流れはイベント駆動モデルに似ている．しかしイベントハンドラとは違い，
これらのハンドラは要求に割り当てられたスレッド上で実行され，ブロック
しうる．

\begin{figure}[htbp]
  \centering
  % Apache: (a) the core invokes modules registered for the stages of request
% processing; (b) several processes, each a boss with a dynamic worker pool.
% Usage: % Apache: (a) the core invokes modules registered for the stages of request
% processing; (b) several processes, each a boss with a dynamic worker pool.
% Usage: % Apache: (a) the core invokes modules registered for the stages of request
% processing; (b) several processes, each a boss with a dynamic worker pool.
% Usage: \input{figures/l06-apache}   (width ~116 mm)
\begin{tikzpicture}[x=1mm, y=1mm, font=\footnotesize\sffamily,
  box/.style={draw, rounded corners=2pt, fill=ln-box, align=center, inner sep=1pt},
  mod/.style={box, fill=white, minimum width=11mm, minimum height=8mm,
    font=\scriptsize\sffamily},
  th/.style={draw, circle, minimum size=5mm, inner sep=0pt,
    font=\scriptsize\sffamily},
  lab/.style={font=\scriptsize\sffamily, text=ln-muted, align=center},
  ttl/.style={font=\footnotesize\sffamily\bfseries, anchor=west},
  >={Stealth[length=1.4mm]}]
  % (a) core and modules
  \node[ttl] at (0,31) {\iflangja{(a) 要求の処理}{(a) request processing}};
  \node[box, fill=ln-accent!15, minimum width=48mm, minimum height=7mm]
    (core) at (24,20) {\iflangja{コア：処理の各段階}{core: stages of processing}};
  \foreach \x/\t in {6/{\iflangja{認証}{authen-\\tication}},
                     18/{\iflangja{ファイル\\提供}{file\\serving}},
                     30/{CGI},
                     42/{\iflangja{ログ}{logging}}} {
    \node[mod] (m\x) at (\x,5) {\t};
    \draw[<->] (\x,16.5) -- (m\x.north);
  }
  \node[lab] at (24,-3) {\iflangja{モジュール（フックに登録）}{modules, registered at hooks}};
  % (b) processes
  \node[ttl] at (56,31) {\iflangja{(b) プロセスとスレッド}{(b) processes and threads}};
  \foreach \px/\n in {56/1,84/2} {
    \draw[rounded corners=2pt] (\px,0) rectangle (\px+25,26);
    \node[th, fill=ln-accent!25] at (\px+12.5,20.5) {B};
    \foreach \i in {0,1,2} {
      \node[th, fill=ln-box] at (\px+5.5+7*\i,12.5) {W};
    }
    \node[th, fill=ln-box] at (\px+5.5,6) {W};
    \node[th, densely dashed] at (\px+12.5,6) {};
    \node[th, densely dashed] at (\px+19.5,6) {};
    \foreach \i in {0,1,2} { \draw[->] (\px+12.5,18) -- (\px+5.5+7*\i,15.2); }
    \node[font=\scriptsize\sffamily, text=ln-muted] at (\px+12.5,1.6)
      {\iflangja{プロセス}{process} \n};
  }
  \node at (113,13) {\dots};
  \node[lab] at (86,-3) {\iflangja{プロセス数もスレッド数も負荷に応じて増減}{processes and thread pools grow and shrink with load}};
\end{tikzpicture}
   (width ~116 mm)
\begin{tikzpicture}[x=1mm, y=1mm, font=\footnotesize\sffamily,
  box/.style={draw, rounded corners=2pt, fill=ln-box, align=center, inner sep=1pt},
  mod/.style={box, fill=white, minimum width=11mm, minimum height=8mm,
    font=\scriptsize\sffamily},
  th/.style={draw, circle, minimum size=5mm, inner sep=0pt,
    font=\scriptsize\sffamily},
  lab/.style={font=\scriptsize\sffamily, text=ln-muted, align=center},
  ttl/.style={font=\footnotesize\sffamily\bfseries, anchor=west},
  >={Stealth[length=1.4mm]}]
  % (a) core and modules
  \node[ttl] at (0,31) {\iflangja{(a) 要求の処理}{(a) request processing}};
  \node[box, fill=ln-accent!15, minimum width=48mm, minimum height=7mm]
    (core) at (24,20) {\iflangja{コア：処理の各段階}{core: stages of processing}};
  \foreach \x/\t in {6/{\iflangja{認証}{authen-\\tication}},
                     18/{\iflangja{ファイル\\提供}{file\\serving}},
                     30/{CGI},
                     42/{\iflangja{ログ}{logging}}} {
    \node[mod] (m\x) at (\x,5) {\t};
    \draw[<->] (\x,16.5) -- (m\x.north);
  }
  \node[lab] at (24,-3) {\iflangja{モジュール（フックに登録）}{modules, registered at hooks}};
  % (b) processes
  \node[ttl] at (56,31) {\iflangja{(b) プロセスとスレッド}{(b) processes and threads}};
  \foreach \px/\n in {56/1,84/2} {
    \draw[rounded corners=2pt] (\px,0) rectangle (\px+25,26);
    \node[th, fill=ln-accent!25] at (\px+12.5,20.5) {B};
    \foreach \i in {0,1,2} {
      \node[th, fill=ln-box] at (\px+5.5+7*\i,12.5) {W};
    }
    \node[th, fill=ln-box] at (\px+5.5,6) {W};
    \node[th, densely dashed] at (\px+12.5,6) {};
    \node[th, densely dashed] at (\px+19.5,6) {};
    \foreach \i in {0,1,2} { \draw[->] (\px+12.5,18) -- (\px+5.5+7*\i,15.2); }
    \node[font=\scriptsize\sffamily, text=ln-muted] at (\px+12.5,1.6)
      {\iflangja{プロセス}{process} \n};
  }
  \node at (113,13) {\dots};
  \node[lab] at (86,-3) {\iflangja{プロセス数もスレッド数も負荷に応じて増減}{processes and thread pools grow and shrink with load}};
\end{tikzpicture}
   (width ~116 mm)
\begin{tikzpicture}[x=1mm, y=1mm, font=\footnotesize\sffamily,
  box/.style={draw, rounded corners=2pt, fill=ln-box, align=center, inner sep=1pt},
  mod/.style={box, fill=white, minimum width=11mm, minimum height=8mm,
    font=\scriptsize\sffamily},
  th/.style={draw, circle, minimum size=5mm, inner sep=0pt,
    font=\scriptsize\sffamily},
  lab/.style={font=\scriptsize\sffamily, text=ln-muted, align=center},
  ttl/.style={font=\footnotesize\sffamily\bfseries, anchor=west},
  >={Stealth[length=1.4mm]}]
  % (a) core and modules
  \node[ttl] at (0,31) {\iflangja{(a) 要求の処理}{(a) request processing}};
  \node[box, fill=ln-accent!15, minimum width=48mm, minimum height=7mm]
    (core) at (24,20) {\iflangja{コア：処理の各段階}{core: stages of processing}};
  \foreach \x/\t in {6/{\iflangja{認証}{authen-\\tication}},
                     18/{\iflangja{ファイル\\提供}{file\\serving}},
                     30/{CGI},
                     42/{\iflangja{ログ}{logging}}} {
    \node[mod] (m\x) at (\x,5) {\t};
    \draw[<->] (\x,16.5) -- (m\x.north);
  }
  \node[lab] at (24,-3) {\iflangja{モジュール（フックに登録）}{modules, registered at hooks}};
  % (b) processes
  \node[ttl] at (56,31) {\iflangja{(b) プロセスとスレッド}{(b) processes and threads}};
  \foreach \px/\n in {56/1,84/2} {
    \draw[rounded corners=2pt] (\px,0) rectangle (\px+25,26);
    \node[th, fill=ln-accent!25] at (\px+12.5,20.5) {B};
    \foreach \i in {0,1,2} {
      \node[th, fill=ln-box] at (\px+5.5+7*\i,12.5) {W};
    }
    \node[th, fill=ln-box] at (\px+5.5,6) {W};
    \node[th, densely dashed] at (\px+12.5,6) {};
    \node[th, densely dashed] at (\px+19.5,6) {};
    \foreach \i in {0,1,2} { \draw[->] (\px+12.5,18) -- (\px+5.5+7*\i,15.2); }
    \node[font=\scriptsize\sffamily, text=ln-muted] at (\px+12.5,1.6)
      {\iflangja{プロセス}{process} \n};
  }
  \node at (113,13) {\dots};
  \node[lab] at (86,-3) {\iflangja{プロセス数もスレッド数も負荷に応じて増減}{processes and thread pools grow and shrink with load}};
\end{tikzpicture}

  \caption{Apache ウェブサーバ．(a) モジュールは，コアが提供する要求処理の
    段階にハンドラを登録する．(b) 複数のプロセスがあり，各プロセスはボス
    （B）とワーカー（W）のプールからなる．破線の円は，負荷が増えたときに
    生成されるワーカーである．}
  \label{fig:l06-apache}
\end{figure}

実行コンテキストは，マルチプロセスモデルとマルチスレッドモデルを
組み合わせている（\cref{fig:l06-apache}b）．Apache は複数のプロセスを実行
し，各プロセスは動的なスレッドプールを持つボス・ワーカーのサーバである．
各プールのスレッド数もプロセス数も，例えばアイドルなワーカーの数や未処理の
接続の数によって，負荷に合わせて調整される．

\begin{note}
  Apache~2 では，骨組み自体が1つのモジュールである．接続を受け付け，
  プロセスとスレッドを管理するのは\emph{マルチプロセッシングモジュール}の
  仕事であり，これが実行コンテキストの構造を選ぶ．\texttt{prefork} は
  プロセスだけを使い，\texttt{worker} はここで述べた混成であり，
  \texttt{event} はアイドルな接続を監視するイベントループを加える．
\end{note}

\needspace{6\baselineskip}
\section{実験の方法論}
\label{sec:l06-methodology}

実験による比較は3つの問いに答えなければならない．どのシステムを比較するか，
どの入力を用いるか，どの指標で測るかである．\source{P2L5，実験の方法論；
Pai ら．}Flash の評価はそのそれぞれを例示している．

\subsection{比較の対象}

Flash は，この章で論じたすべてのアーキテクチャのサーバと比較された
（\cref{tab:l06-comparison-points}）．MP，MT，SPED の各サーバは Flash の
コードから作られており，アーキテクチャだけが Flash と異なり，最適化は共通
である．したがってそれらの間の差は，アーキテクチャに帰することができる．
2つの独立したサーバは，当時使われていたサーバの中で Flash がどの位置に
あるかを示す．

\begin{table}[htbp]
  \centering
  \caption{Flash の評価で比較されたサーバ．}
  \label{tab:l06-comparison-points}
  \small
  \begin{tabular}{@{}l>{\raggedright\arraybackslash}p{0.72\linewidth}@{}}
    \toprule
    サーバ & アーキテクチャ \\
    \midrule
    MP     & マルチプロセス．各プロセスは単一スレッドで，一度に1つの要求を
             扱う \\
    MT     & マルチスレッド，ボス・ワーカー \\
    SPED   & 単一プロセスイベント駆動，ヘルパなし \\
    Zeus   & バージョン 1.30，商用の SPED サーバ．トレースに基づく実験では
             2つのプロセスで実行 \\
    Apache & バージョン 1.3.1，マルチプロセス \\
    Flash  & AMPED \\
    \bottomrule
  \end{tabular}
\end{table}

\subsection{ワークロード}

\begin{definition}[ワークロード]
  実験の\term[わーくろーど]{ワークロード}[workload]とは，試験対象の
  システムに与える入力の集合である．ウェブサーバであれば，到着時刻と要求
  されるファイルを伴う要求の列である．
\end{definition}

ワークロードは現実的であるべきで，例えば実際のサーバが受けるウェブページ
要求の時間分布を再現すべきである．また制御可能かつ再現可能であって，どの
サーバも同じ入力で測定されるようにすべきである．
\term[とれーす]{トレース}[trace]，すなわち実際のシステムで観測された要求の
記録を実験で再生するものは，この両方の要件を満たす．
\term[ごうせいわーくろーど]{合成ワークロード}[synthetic workload]は，選んだ
パラメータからプログラムが生成するものであり，現実性は劣るが，各パラメータ
を系統的に変化させられる．Flash は，ライス大学のウェブサーバで収集された
2つのトレースで評価された．大きなファイル集合を持つ計算機科学科の
\emph{CS トレース}と，学生の個人ページを提供するサーバの，より小さな集合を
持つ \emph{Owlnet トレース}である．合成ワークロードは，大きさを変えながら
1つのファイルを繰り返し要求するものであった．

\subsection{指標}

ウェブサーバは，届けるデータと応じるクライアントによって，2つの指標で
測られる．
\begin{equation}
  \begin{aligned}
    \text{帯域幅} &= \frac{\text{転送した総バイト数}}{\text{総時間}} ,\\[1mm]
    \text{接続レート} &= \frac{\text{クライアント接続の総数}}{\text{総時間}}
  \end{aligned}
  \label{eq:l06-bandwidth}
\end{equation}
どちらもファイルサイズの関数として評価された．ファイルが大きいほど，サーバ
は接続の固定コストをより多くのバイトに分散でき，帯域幅は上がる．しかし同時
に接続当たりの作業も増えるので，接続レートは下がる．これは簡単なモデルで
定量的に表せる．大きさ $S$ のファイルを提供するのに，接続当たり固定の時間
$c$ とバイト当たり $b$ がかかるとすれば，一度に1つの接続を扱うサーバは
\begin{equation}
  R(S) = \frac{1}{c + b\,S} ,
  \qquad
  B(S) = S \cdot R(S) = \frac{S}{c + b\,S} \;<\; \frac{1}{b}
  \label{eq:l06-amortize}
\end{equation}
を達成する．ここで $R$ は接続レート，$B$ は帯域幅である．%
\tip{リトルの法則が両者を結ぶ．レート $R$，平均応答時間 $T$ なら同時に
開いている接続は $R\,T$ 個であり，毎秒 \num{1000} 接続で
\SI{100}{\milli\second} なら \num{100} 個である．}

\begin{example}
  $c = \SI{1}{\milli\second}$，キロバイト当たり
  $b = \SI{0.02}{\milli\second}$ とすると，\SI{10}{\kilo\byte} のファイル
  には \SI{1.2}{\milli\second} かかり，$R \approx 833$ 接続毎秒，
  $B \approx \SI[per-mode=symbol]{8.3}{\mega\byte\per\second}$ となる．
  \SI{100}{\kilo\byte} のファイルには \SI{3}{\milli\second} かかり，
  $R \approx 333$ 接続毎秒だが，
  $B \approx \SI[per-mode=symbol]{33.3}{\mega\byte\per\second}$ となる．
  ファイルが大きくなるにつれて $B$ は
  $1/b = \SI[per-mode=symbol]{50}{\mega\byte\per\second}$ に近づき，一方
  $R$ は下がり続ける．
\end{example}

\needspace{6\baselineskip}
\section{実験結果}
\label{sec:l06-results}

評価の結果は何よりも，要求されたデータがメモリ上にあるかどうかに依存する．
\source{P2L5，実験結果，性能結果のまとめ；Pai ら．}\cref{tab:l06-results}は
以下で述べる知見をまとめたものである．

\begin{table}[htbp]
  \centering
  \caption{Flash の評価における性能結果のまとめ．}
  \label{tab:l06-results}
  \small
  \begin{tabular}{@{}>{\raggedright\arraybackslash}p{0.21\linewidth}>{\raggedright\arraybackslash}p{0.24\linewidth}>{\raggedright\arraybackslash}p{0.47\linewidth}@{}}
    \toprule
    ワークロード & 順位 & 理由 \\
    \midrule
    データがキャッシュ上 & SPED $>$ Flash & Flash はメモリ上にあるかの判定
      を不要に行う \\
    & SPED，Flash $>$ MT，MP & 同期とコンテキストスイッチのオーバーヘッド \\
    \midrule
    ディスクを多用 & Flash $>$ SPED & SPED は非同期I/O を欠きブロックする \\
    & Flash $>$ MT，MP & Flash はメモリ効率が高く，コンテキストスイッチも
      少ない \\
    \bottomrule
  \end{tabular}
\end{table}

\begin{description}
  \item[単一ファイル（最良の場合）] 合成ワークロードでは，すべての要求が
    同じファイルに対するもので，そのファイルは常にキャッシュされている．
    帯域幅，すなわち $n$ 個の要求とファイルサイズの積を時間で割った値が，
    \SI{200}{\kilo\byte} までのファイルサイズについて測定された．%
    \margin{サーバマシンは \SI{333}{\mega\hertz} の Pentium~II で，メモリ
    \SI{128}{\mega\byte}，\SI[per-mode=symbol]{100}{\mega\bit\per\second} の
    インタフェースを複数持ち，Solaris~2.6 と FreeBSD~2.2.6 で動作した
    \cite{pai1999}．}%
    どのサーバも似た傾向を示す．最も性能が高いのは SPED である．Flash はそれに
    迫るが，ここでは決して必要にならない \texttt{mincore} の判定を要求
    ごとに行う分を支払う．Zeus はいくつかのファイルサイズで異常を示し，
    Pai らはこれを上で述べた整列の問題に帰している．MT と MP は同期と
    コンテキストスイッチのために遅く，最適化を欠く Apache が最も遅い．
  \item[Owlnet トレース] データ集合は小さく，その大半がキャッシュに収まる
    ので，傾向は最良の場合に似る．しかし時折，ディスクを必要とする要求が
    ある．そのとき SPED は完全にブロックしてしまうが，Flash はヘルパの
    おかげで他の要求に応じ続けられる．
  \item[CS トレース] データ集合はより大きく，多くの要求がディスク I/O を
    必要とする．Flash のコードから作られたサーバの中では，非同期I/O を
    欠く SPED が最も性能が低い．MT は MP より良い．メモリ使用量が小さい分
    だけファイルのキャッシュに使えるメモリが増えてディスク I/O が減り，
    スレッド間で状態を共有する方が，プロセスに必要なプロセス間通信より
    安価だからである．最も性能が高いのは Flash である．メモリ使用量が
    小さいのでキャッシュに使えるメモリが増え，ブロッキング I/O を必要と
    する要求が減り，しかも同期を必要としない．Apache はここでも，測定
    されたすべてのサーバの中で最も遅い．%
    \margin{測定したワークロード全体で，Flash は「既存のウェブサーバの
    性能に並ぶか，最大で 50\,\% 上回る」\cite{pai1999}．}
  \item[最適化] パス名，応答ヘッダ，対応付けたファイルの各キャッシュを
    さまざまに組み合わせて Flash を測定すると，どの最適化も相当に寄与して
    いることが分かる．Apache もこれらから恩恵を受けたはずである．
\end{description}

\begin{keypoint}
  データがメモリ上にあるときはイベント駆動のサーバが勝ち，Flash の余分な
  判定がわずかな代償となる．ワークロードがディスクを必要とするときは，
  Flash のヘルパがブロックを避け，メモリ使用量が小さいのでキャッシュに
  使えるメモリが増え，ディスク I/O が減る．Flash は個々のワークロードで
  最良というより，ワークロード全体にわたって良い性能を示す．
\end{keypoint}

\needspace{6\baselineskip}
\section{実験の設計と実施}
\label{sec:l06-design}

実験を行う価値があるのは，それが関連性を持つ場合に限られる．すなわち，他者
が信じ，関心を持つ解決策についての主張を支えるものでなければならない．
\source{P2L5，実験の設計に関する助言，実験の実施に関する助言．}

\subsection{関連性と指標}

関連性は関係者に依存する．ウェブサーバの場合，クライアントは応答時間を
気にし，サーバの運用者はスループットと費用を気にする．%
\sidenote{クライアントに平均ではなくパーセンタイルを示すのは，最も遅い応答
が全体を決めるからである．1つの利用者要求が \num{100} 台のサーバの応答を
待ち，\num{100} 回に1回が1秒を超えるなら，その要求が遅くなる確率は
$1 - 0.99^{100} \approx 63\,\%$ である．Dean と Barroso はこれを
\emph{tail at scale} と呼ぶ \cite{dean2013}．}%
応答時間とスループットの両方を改善する新しい設計は，明らかに価値がある．一方を改善し他方を悪化
させない設計も，やはり改善である．他方を犠牲にして一方を改善する設計も，
悪化した指標が許容範囲に留まるならば，一部の関係者には有用でありうる．要求
レートが増えても応答時間を変えないことも，目標になりうる．このように目標が，
どの指標を測るか，実験をどう構成するかを決める．

指標はそれに応じて選ぶべきである．標準的な指標はより広い読者に届く．読者が
その解釈と比較の仕方を知っているからである．同じ理由で，実験ではしばしば
\term[べんちまーく]{ベンチマーク}[benchmark]，すなわちコミュニティが合意
した，実施方法と結果の報告の規則を伴う標準的なワークロードが使われる．
それに加えて，指標は，誰が結果に関心を持つのか，何に関心を持つのか，なぜ
かという問いに答えるものであるべきである．クライアントの性能であれば応答
時間やタイムアウトした要求の数を，運用者の費用であればスループットと費用
を求めることになる．

\subsection{構成空間}

システムの振る舞いは多くの要因に依存し，それらが合わさってその構成空間を
なす．
\begin{itemize}
  \item \textbf{システム資源}：CPU やメモリのようなハードウェアと，
    スレッド数やキューの大きさのようなソフトウェア．
  \item \textbf{ワークロード}：ウェブサーバであれば，要求レート，並行する
    要求の数，ファイルサイズ，アクセスパターン．
\end{itemize}
空間全体を探索することはできないので，実験では構成パラメータの部分集合を
選び，変化させる各パラメータについて値の範囲を決め，関連性のある
ワークロードを用いる．最良の場合と最悪の場合のシナリオは含めるべきである．
要因の組み合わせは有用なものであるべきで，多くの組み合わせは同じ点を
繰り返すだけである．

比較を律する規則が2つある．第1に，\emph{同じもの同士}を比較することである．
結果を比較する2つの構成は，調べている要因においてのみ異なっていなければ
ならない．あるサーバを大きなマシン上の軽いワークロードで測り，別のサーバを
小さなマシン上の重いワークロードで測ったのでは，サーバについて何も結論
できない．差がワークロードから来るのか，資源から来るのか，サーバから来るの
か分からないからである．第2に，\term[べーすらいん]{ベースライン}[baseline]
と比較することである．これは性能が理解されている参照システムであり，最先端
のもの，最も一般的な手法，あるいは理想的な最良・最悪の場合のシナリオなどで
ある．ある設計が優れているという主張は，それがベースラインより優れていると
いう主張である．%
\caution{ベースラインは新しい設計と同じだけ丁寧に調整しなければならない．
既定の設定のまま比較することが，改善を過大に見せる最もありがちなやり方で
ある．}

\subsection{実験の実施}

実験を設計したら，各試験事例を何度も実行し，実行にわたって指標を，例えば
平均とばらつきとして計算し，結果を，主張が明確に伝わる形式---典型的には
グラフ---で表す．最後に，結論を導かなければならない．実験の目標に照らして
解釈されない測定値は，どんな主張も支えない．

\begin{example}
  あるチームが，新しい AMPED のサーバが写真共有サイトの応答時間の 95
  パーセンタイルを，スループットを下げずに短縮することを示したいとする．
  ハードウェアは1台のサーバマシンに固定し，ワークロードはそのサイトの
  トレースを，毎秒 1\,000，2\,000，4\,000，8\,000，16\,000 要求と要求
  レートを上げながら再生したものとする．すべてのファイルがメモリ上にある
  最良の場合と，データ集合がメモリの4倍ある最悪の場合とが範囲を挟む．
  ベースラインは現在使われているマルチスレッドのサーバである．2つのサーバ，
  5つのレート，2つの場合，それぞれ5回の繰り返しで，この研究は
  $2 \cdot 5 \cdot 2 \cdot 5 = 100$ 回の実行からなり，各構成について
  パーセンタイルとスループットの平均を，そのばらつきとともに報告する．
\end{example}

\needspace{11\baselineskip}
\begin{keypoint}[章のまとめ]
  どの設計が優れているかは，指標とワークロードに依存する．並行サーバは
  プロセス（MP），スレッド（MT），あるいは1つのイベント駆動のスレッド
  （SPED）を使える．非同期I/O が利用できないとき，ヘルパはブロッキング操作
  をイベントに変える（AMPED，AMTED）．AMPED のサーバである Flash は，
  データがキャッシュされているときは SPED に迫り，ワークロードがディスクを
  必要とするときは他の選択肢に並ぶかそれを上回る．Apache はプロセスと
  動的なスレッドプールを組み合わせている．健全な実験は，関連性のある指標と
  ワークロードを選び，調べている要因において異なる構成同士を比較し，
  ベースラインと比較し，実行を繰り返し，結論を導く．
\end{keypoint}

\end{document}
